%# -*- coding:utf-8 -*-


\chapter{三角}


\section{正弦、余弦、正切、余切}


\subsection{锐角的正弦、余弦、正切、余切}
\zsd 锐角的正弦、余弦、正切、余切

\begin{enumerate}
\item 填空\\
\begin{minipage}{\textwidth}
\centering
\includegraphics[scale=0.9]{./bodyxb/src/main42xbb2ch06ruijiandesanjiaobi01t1.png}
\end{minipage}
\begin{enumerate}[(1)]
\item $\sin{A}=$\blkda{}, $\cos{A}=$\blkda{}, $\tan{A}=$\blkda{}, $\cot{A}$=\blk. 
\item $\sin^2A+\cos^2A=$\blkda{}. 
\item $\sin(\ang{90}-A)=$\blkda{}, $\cos(\ang{90}-A)=$\blkda{}, \\
$\tan(\ang{90}-A)=$\blkda{}, $\cot(\ang{90}-A)=$\blkda{}. 
\item 
\begin{minipage}[t]{\textwidth}
\begin{tblr}{hlines,vlines,column{1}={2cm,c},column{2,3,4,5}={2.5cm,c},row{1}={blue!10}}
角度$\alpha$ & $\sin\alpha$ & $\cos\alpha$ & $\tan\alpha$ & $\cot\alpha$\\
$\ang{30}$ & & & & \\
$\ang{45}$ & & & & \\
$\ang{60}$ & & & & \\
\end{tblr}
\end{minipage}
\end{enumerate}

\item 如图, 点$O$为$\triangle{ABC}$边$AC$的中点, 连接$BO$并延长到点$D$, 连接$AD$、$CD$, 
若$BD=12$, $AC=8$, $\angle{AOD}=\ang{120}$, 则四边形$ABCD$的面积为\blkda{$96\sqrt{3}$}. 

\item 如图, 矩形纸片$ABCD$, $AB=4$, $BC=3$, 点$P$在$BC$边上 , 将$\triangle{CDP}$沿$DP$折叠, 
点$C$落在点$E$处, $PE$、$DE$分别交$AB$于点$O$、$F$, $OP=OF$, 则$\cos\angle{ADF}$的值为\blkda{3}. 
\begin{jieda}
$BE$在$BC$上的投影长度为$\dfrac{1}{2}$. 
\end{jieda}

\item 如图, 在矩形$ABCD$中, $BC=2$, $AE\bot BD$, 垂足为$E$, $\angle{BAE}=\ang{30}$, 则$\tan\angle{DEC}$的值是\blkda{}. 

\begin{figure}[!h]
\centering
\begin{minipage}[b]{0.3\textwidth}
\centering
\includegraphics[totalheight=6\baselineskip]{./bodyxb/src/main42xbb2ch06ruijiandesanjiaobi01t2.png}
\figcaption*{第2题}
\end{minipage}
\qquad
\begin{minipage}[b]{0.3\textwidth}
\centering
\includegraphics[totalheight=6\baselineskip]{./bodyxb/src/main42xbb2ch06ruijiandesanjiaobi01t3.png}
\figcaption*{第3题}
\end{minipage}
\begin{minipage}[b]{0.3\textwidth}
\centering
\includegraphics[totalheight=5.5\baselineskip]{./bodyxb/src/main42xbb2ch06ruijiandesanjiaobi01t4.png}
\figcaption*{第4题}
\end{minipage}
\end{figure}



\item 如图, 从点$A$看一山坡上的电线杆$PQ$, 观测点$P$的仰角是$\ang{45}$, 向前走6\si{m}到达$B$点, 
测得顶端点$P$和杆底端点$Q$的仰角分别是$\ang{60}$和$\ang{30}$, 则该电线杆$PQ$的高度\blkda{}. 
\begin{jieda}

\end{jieda}

\item ``奔跑吧, 兄弟!''节目组, 预设计一个新的游戏: ``奔跑''路线需经$A$、$B$、$C$、$D$四地. 
如图, 其中$A$、$B$、$C$三地在同一直线上, $D$地在$A$地北偏东$\ang{30}$方向在$C$地北偏西$\ang{45}$方向. 
$C$地在$A$地北偏东$\ang{75}$方向, 且$BD=BC=30\si{m}$, 则从$A$地到$D$地的距离是\blkda{$15\sqrt{6}$}. 
\begin{jieda}

\end{jieda}

\item 如图, 两根竹竿$AB$和$AD$斜靠在墙$CE$上, 量得$\angle{BAC}=\ang{60}$, $\angle{DAC}=\ang{70}$, 则竹竿$AB$与$AD$的长度之比为\dotfill(\kh{B})
\xx{$2\sin\ang{70}$}{$2\cos\ang{70}$}{$2\tan\ang{70}$}{$\dfrac{2}{\tan\ang{70}}$}

\begin{figure}[!h]
\centering
\begin{minipage}[b]{0.3\textwidth}
\centering
\includegraphics[totalheight=6\baselineskip]{./bodyxb/src/main42xbb2ch06ruijiandesanjiaobi01t5.png}
\figcaption*{第5题}
\end{minipage}
\qquad
\begin{minipage}[b]{0.3\textwidth}
\centering
\includegraphics[totalheight=6\baselineskip]{./bodyxb/src/main42xbb2ch06ruijiandesanjiaobi01t6.png}
\figcaption*{第6题}
\end{minipage}
\qquad
\begin{minipage}[b]{0.3\textwidth}
\centering
\includegraphics[totalheight=6\baselineskip]{./bodyxb/src/main42xbb2ch06ruijiandesanjiaobi01t7.png}
\figcaption*{第7题}
\end{minipage}
\end{figure}
\begin{jieda}
$AC=AB\cos\ang{60}=AD\cos\ang{70}$. 
\end{jieda}


\item 如图, 基灯塔$AB$建在陡峭的山坡上, 该山坡的坡度$i=1:0.75$. 小明为了测得灯塔的高度, 他首先测得$BC=20\si{m}$, 
然后在$C$处水平向前走了34\si{m}到达一建筑物底部$E$处, 他在该建筑物顶端$F$处测得灯塔顶端A的仰角为$\ang{43}$. 
若该建筑物$EF=20\si{m}$, 则灯塔$AB$的高度约为(精确到0.1\si{m})\dotfill(\kh{})
\xx{46.7\si{m}}{46.8\si{m}}{53.5\si{m}}{67.8\si{m}}
\begin{jieda}
$AB=AG+EF=GF\cdot\tan\ang{43}+EF=(12+34)\tan\ang{43}+20=\doteq 62.9$. \\
\ps 坡度的定义是$\dfrac{\Delta V}{\Delta H}=\tan\theta$. 
\end{jieda}

\item 如图, 一艘轮船从位于灯塔$C$的北偏东$\ang{60}$方向, 距离灯塔$60n~\si{mile}$的小岛$A$出发, 沿正南方向航行一段时间后, 
到达位于灯塔$C$的南偏东$\ang{45}$方向上的$B$处, 这时轮船$B$与小岛$A$的距离是\dotfill(\kh{})
\xx{$30\sqrt{3}n~\si{mile}$}
{$60n~\si{mile}$}
{$120n~\si{mile}$}
{$(30+30\sqrt{3})n~\si{mile}$}
\begin{jieda}

\end{jieda}

\item 如图, 在$x$轴的上方, 直角$\angle{BOA}$绕原点$O$按顺时针方向旋转. 若$\angle{BOA}$的两边分别与函数$y=-\dfrac{1}{x}$、$y=\dfrac{2}{x}$
的图像交于$B$、$A$两点, 则$\angle{OAB}$大小的变化趋势为\dotfill(\kh{D})
\xx{逐渐变小}{逐渐变大}{时大时小}{保持不变}

\begin{figure}[!ht]
\centering
\begin{minipage}[b]{0.3\textwidth}
\centering
\includegraphics[totalheight=6\baselineskip]{./bodyxb/src/main42xbb2ch06ruijiandesanjiaobi01t8.png}
\figcaption*{第2题}
\end{minipage}
\qquad
\begin{minipage}[b]{0.3\textwidth}
\centering
\includegraphics[totalheight=6\baselineskip]{./bodyxb/src/main42xbb2ch06ruijiandesanjiaobi01t9.png}
\figcaption*{第9题}
\end{minipage}
\qquad
\begin{minipage}[b]{0.3\textwidth}
\centering
\includegraphics[totalheight=6\baselineskip]{./bodyxb/src/main42xbb2ch06ruijiandesanjiaobi01t10.png}
\figcaption*{第10题}
\end{minipage}
\end{figure}
\begin{jieda}
$\tan\angle{OAB}=\dfrac{OB}{OA}=\dfrac{\sqrt{1+\frac{1}{k^2}}\cdot\sqrt{k}}{\sqrt{1+k^2}\cdot\sqrt{\frac{2}{k}}}=\dfrac{1}{\sqrt{2}}$. \\
其中$k$是直线$OA$的斜率. 
\end{jieda}

\end{enumerate}





\subsection{任意角及其度量(1)}
\zsd 角度制、任意角、象限角、终边相同的角 \\
本节: 角的顶点在坐标原点, 始边与$x$轴正半轴重合. 

\begin{enumerate}
\item 判断下列命题是否正确, 并说明理由: 
\begin{enumerate}
\item 小于$90^{\circ}$的角是锐角; 
\item 第一象限角小于第二象限角; 
\item 终边相同的角一定相等; 
\item 相等的角终边一定相同; 
\item 若$90^{\circ}<\alpha<180^{\circ}$, 则$\alpha$是第二象限角. 
\end{enumerate}
\begin{jieda}
\begin{enumerate}[(1)]
\item $-\ang{45}$角不是锐角, 故假. 
\item $\ang{30}>\ang{120}-\ang{360}=-\ang{240}$, 故假. 
\item $\alpha=\ang{30}$, $\beta=\ang{30}+\ang{360}$, 故假. 
\item 真. 如果有破绽的话, 破绽在于始边是否重合. 
\item 真. 
\end{enumerate}
\end{jieda}

\item $-2021^{\circ}$是第\blkda{II}象限角. 
\begin{jieda}
$-2021=-5\cdot360-221=-6\cdot 360+139$, 故II. 
\end{jieda}


\item 若时钟的分针转了2小时15分, 则所转的角度是\blkda{$-\ang{810}$}. 
\begin{jieda}
$s=vt=360\cdot 2.25=810$, 就是这个东西. 
\end{jieda}


\item 如图所示, 终边是射线$OA$的角的集合是\blkda[4]{$\{k\cdot\ang{360}+\ang{60}\mid k\zz\}$}; 
终边是射线$OB$的角的集合是\blkda[4]{$\{k\cdot\ang{360}+\ang{225}\mid k\zz\}$}. 


\item 角的顶点在坐标系的原点, 始边与$x$轴的正半轴重合, 那么终边在下列位置的角的集合分别是: 
\begin{enumerate}[(1)]
\item $x$轴正半轴: \blkda[5]{$\{\theta\mid \theta=k\cdot\ang{360},k\zz\}$}; 
\item $x$轴: \blkda[5]{$\{\theta\mid \theta=k\cdot\ang{180},k\zz\}$}; 
\item $y$轴负半轴: \blkda[5]{$\{\theta\mid \theta=k\cdot\ang{360}-\ang{90},k\zz\}$}; 
\item $y$轴: \blkda[5]{$\{\theta\mid \theta=k\cdot\ang{180}+\ang{90},k\zz\}$}; 
% \item 坐标轴: \blkda{}; 
\item 直线$y=-x$: \blkda[5]{$\{\theta\mid \theta=k\cdot\ang{180}+\ang{135},k\zz\}$}; 
\item 直线$y=\pm{x}$: \blkda[5]{$\{\theta\mid \theta=k\cdot\ang{90}+\ang{45},k\zz\}$}. 
\end{enumerate}
\begin{jieda}
以(6)为例: 每逆时针旋转$\ang{90}$, 恰好到达一条射线. 于是从角的角度看, 遍历了所有的角. 
相当于有了一个起始值$\ang{45}$, 每次加上或减去$\ang{90}$, 就有了所有角. 
\end{jieda}


\item 已知$\alpha$与$\beta$都是锐角, $\alpha+\beta$的终边与$-280^{\circ}$的终边相同; 
$\alpha-\beta$的终边与$-670^{\circ}$的终边相同, 则$\alpha$=\blkda{$\ang{65}$}, $\beta$=\blkda{$\ang{15}$}. 
\begin{jieda}
$\begin{cases}
\alpha+\beta=m\cdot\ang{360}-\ang{280}\\
\alpha-\beta=n\cdot\ang{360}-\ang{670}
\end{cases}$, 
故$\alpha=(m+n)\cdot\ang{180}-\ang{475}$, $\beta=(m-n)\cdot\ang{180}+\ang{195}$. \\
又$\alpha,\beta$是锐角, 故$\alpha=\ang{65}$, $\beta=\ang{15}$. 
\end{jieda}


\item 已知圆周上有一点$P$(初始位置在$x$轴正半轴上)依逆时针方向作匀速圆周运动. 已知点$P$每分钟转过$\theta$角($\ang{0}<\theta\leq\ang{180})$, 
2分钟到达第三象限, 15分钟回到原来的位置, 则角$\theta$等于\blkda[2]{$\ang{96}$或$\ang{120}$}. 
\begin{jieda}
$\begin{cases}
\ang{0}<\theta<\ang{180}\\
\ang{180}<2\theta<\ang{270}\\
15\theta=k\cdot\ang{360}
\end{cases}$, 
将第三式代入前两式得到$k$的范围, 结合$k\zz$得$k=4,5$, 故$\theta=\ang{96}$或$\ang{120}$. 
\end{jieda}


\item (选做)某天下午小明约同学去踢球, 玩了3个多小时, 离家时他看了时钟是2点多钟, 回家时又看了钟, 发现时针与分针正好交换了位置, 
则小明离开家的时间总长是\blkda{222分钟}(精确到分钟). 
\begin{jieda}
角度不同, 过程的多少有别. \\
时针转了$\theta$角, 则分针转了$12\theta$, 且分针共超过时针3次, 
故$12\theta=4\cdot\ang{360}-\theta$. 
\end{jieda}

\item 角$\alpha$是第四象限角, 则角$\ang{180}-\alpha$是\dotfill(\kh{C})
\xx{第一象限角}
{第二象限角}
{第三象限角}
{第四象限角}
\begin{jieda}
另外一个角度: $\dfrac{\alpha+(\ang{180}-\alpha)}{2}=\ang{90}$, \\
$\ang{90}$对应于$y$正半轴, 故两者关于$y$轴正半轴对称. \\
实际上, 关于$y$轴正半轴与关于$y$轴是等效的. 
\end{jieda}


\item 若$\alpha$与$\beta$是终边相同的角, 则$\alpha-\beta$的终边一定在\dotfill(\kh{A})
\xx{$x$轴的非负半轴上}
{$x$轴的非正半轴上}
{$y$轴的非正半轴上}
{$y$轴的非负半轴上}



\item 设$k\zz$, 终边相同的一组角是\dotfill(\kh{C})
\xx{$k\cdot 90^{\circ}$与$k\cdot 180^{\circ}+90^{\circ}$}
{$k\cdot 180^{\circ}+60^{\circ}$与$k\cdot 60^{\circ}$}
{$(2k-1)\cdot 180^{\circ}$与$(4k\pm 1)\cdot 180^{\circ}$}
{$k\cdot 180^{\circ}+30^{\circ}$与$k\cdot 360^{\circ}-150^{\circ}$}



\item 分别写出第一、二、三、四象限角的集合. 
\begin{jieda}
一: $\{\theta\mid\theta=k\cdot\ang{360}+\alpha,\ang{0}\alpha<\ang{90},k\zz\}$; \\
二: $\{\theta\mid\theta=k\cdot\ang{360}+\alpha,\ang{90}\alpha<\ang{180},k\zz\}$; \\
三: $\{\theta\mid\theta=k\cdot\ang{360}+\alpha,\ang{180}\alpha<\ang{270},k\zz\}$; \\
四: $\{\theta\mid\theta=k\cdot\ang{360}+\alpha,\ang{270}\alpha<\ang{360},k\zz\}$. 
\end{jieda}
\end{enumerate}




\subsection{任意角及其度量(2)}
\zsd 弧度制、扇形及其弧长、面积公式

\begin{enumerate}
\item $\dfrac{\pi}{5}$=\blkda{36}度; $-72^{\circ}$=\blkda{$-\dfrac{2\pi}{5}$}弧度. 

\item 在与2021弧度角的终边相同的角中, 绝对值最小的角是\blkda[3]{$2021-644\pi$}. 
\begin{jieda}
$\abs{2021+2k\pi}_{\min}=\abs{2021-322\cdot2\pi}$. 
\end{jieda}


\item 若角$\alpha$和角$\beta$的和是$1$弧度, 差为$1^{\circ}$, 则$\alpha$和$\beta$的弧度数分别为\blkda{$\dfrac{1}{2}\pm\dfrac{1}{360}$}. 

\item 若中心角为$1^{\circ}$的扇形, 它的弧长为$2\pi$, 则该扇形所在圆的半径为\blkda{360}. 
\begin{jieda}
$r=\dfrac{l}{\alpha}=\dfrac{2\pi}{\frac{\pi}{180}}=360$. 
\end{jieda}


\item 若扇形的圆心角是$120^{\circ}$, 则该扇形面积与其内切圆面积的比值是\blkda{$\dfrac{7+4\sqrt{3}}{9}$}. 
\begin{jieda}
相切找切点, 于是结合简单的几何知识易得$R=r+\dfrac{2r}{\sqrt{3}}$. 
\end{jieda}


\item (选做题)若在直径为$10\si{cm}$的滑轮上有一条弦, 其长为$6\si{cm}$, 且$P$为弦的中点, 滑轮以每秒$5$弧度的角速度旋转, 则经过$5$秒后, $P$点转过的弧长是\blkda{$100\si{cm}$}. 


\item 已知地球赤道的半径是$6370\si{km}$, 所以赤道上${1}'$的弧长是(精确到$0.01\si{km})$\dotfill(\kh{A})
\xx{$1.85\si{km}$}{$1.88\si{km}$}{$111.18\si{km}$}{112.36\si{km}}


\item 若一条弦长等于半径的$\dfrac{1}{2}$, 则此弦所对圆心角\dotfill(\kh{D})
\xx{等于$\dfrac{\pi}{6}$弧度}
{等于$\dfrac{\pi}{3}$弧度}
{等于$\dfrac{1}{2}$弧度}
{以上都不对}
\begin{jieda}
$\sin\frac{\theta}{2}=\frac{1}{4}$. 
\end{jieda}


\item 若扇形的周长是$16$, 圆心角是$2$弧度, 则扇形面积是\dotfill(\kh{C})
\xx{$16\pi$}{$32\pi$}{$16$}{$32$}


\item 如果$1$弧度的圆心角所对的弦长为$2$, 那么这个圆心角所对的弧长是\dotfill(\kh{C})
\xx{$\sin\dfrac{1}{2}$}
{$\dfrac{\pi}{6}$}
{$\dfrac{1}{\sin\frac{1}{2}}$}
{$2\sin\dfrac{1}{2}$}
\begin{jieda}
``弧''与``弦''; 弦取中点. 
\end{jieda}


\item 一扇形周长是\SI{32}{cm}, 扇形的圆心角为多少弧度时, 这个扇形的面积最大? 最大面积是多少? 
\begin{jieda}
$l+2r=32$, $S=\frac{1}{2}lr=\frac{1}{4}\cdot l\cdot2r\leq\frac{1}{4}(\frac{l+2r}{2})^2=64$. \\
$\alpha=\frac{l}{r}=\SI{2}{rad}$, $S_{\min}=\SI{64}{cm^2}$. 
\end{jieda}


\item
\begin{enumerate}[(1)]
\item 若角$\alpha$的终边过点$(2^x,\log_{\frac{1}{2}}2)$, $x\zr$, 则$\dfrac{\alpha}{2}$终边在第\blkda{II、IV}象限. 

\item 设$\alpha$是第二象限角, 则$\dfrac{\alpha}{3}$的终边不在\dotfill(\kh{C})
\xx{第一象限}
{第二象限}
{第三象限}
{第四象限}

\item 已知象限角$\alpha$, 如何快速判断角$\dfrac{\alpha}{n}$所在的象限? 
\end{enumerate}
\begin{jieda}
先看图: \\
\begin{minipage}{\textwidth}
\centering
\includegraphics[scale=0.8]{./bodyxb/src/main42chap05renyijiaodesanjiaobi01.pdf}
\figcaption{半角、倍角与象限}
\end{minipage}
该图的信息有点多: 
\begin{itemize}
\item 由$\alpha$所在的象限, 可得$\dfrac{\alpha}{2}$所在象限, 且只能是象限中的一部分. 
\item 由$\alpha$所在的象限, 可得$2\alpha$所在象限. 
\item 可推广到$\dfrac{\alpha}{n}(n\zn)$的情形. 
\end{itemize}
\begin{enumerate}[(1)]
\item $2^x>0, \log_{\frac{1}{2}}2=-1$, 故$\alpha$是第四象限角, 于是$\dfrac{\alpha}{2}$是第II,IV象限角. 

\item 每个象限等分为3份, 然后逆时针依次标上1234(周期). 

\item 每个象限等分为$n$份, 然后逆时针依次标上1234(周期). 
\end{enumerate}
\end{jieda}

\end{enumerate}






\subsection{任意角的正弦、余弦、正切、余切(1)}
\zsd 任意角的正弦、余弦、正切、余切及其符号

\begin{enumerate}

\item 已知角$\alpha$的终边经过下列各点, 求角$\alpha$的正弦、余弦、正切、余切. 
\begin{enumerate}[(1)]
\item $A(-8,-6)$,     则$\sin\alpha$=\blkda{$-\frac{3}{5}$}, $\cos\alpha$=\blkda{$-\frac{4}{5}$}, $\tan\alpha$=\blkda{$\frac{3}{4}$}, $\cot\alpha$=\blkda{$\frac{4}{3}$}; 
\item $B(\sqrt3,-1)$, 则$\sin\alpha$=\blkda{$-\frac{1}{2}$}, $\cos\alpha$=\blkda{$\frac{\sqrt{3}}{2}$}, $\tan\alpha$=\blkda{$-\frac{\sqrt{3}}{3}$}, $\cot\alpha$=\blkda{$-\sqrt{3}$}. 
\end{enumerate}

\item 已知角$\theta$的终边上有一点$P(4,y)$, 且$\sin\theta=-\dfrac{2\sqrt{5}}{5}$, 则$y$=\blkda{$-8$}. 

\item 已知角$\alpha$的终边与射线$5x+12y=0(x\leq 0)$的图像重合, $\cos\alpha+\cot\alpha-\tan\alpha$的值为\blkda{$-\dfrac{2267}{780}$}. 
\begin{jieda}
$5x+12y=0(x\leq 0)$上取点$(-12,5)$, 
则$E=\dfrac{-12}{13}+\dfrac{-12}{5}-\dfrac{5}{-12}=-\dfrac{2267}{780}$. 
\end{jieda}

\item 已知集合$A=\left\{y\mid y=\dfrac{\abs{\sin{x}}}{\sin{x}}+\dfrac{\cos{x}}{\abs{\cos{x}}}+\dfrac{\abs{\tan{x}}}{\tan{x}}\right\}$, 
用列举法表示集合$A$是\blkda{$\{-1,3\}$}. 
\begin{jieda}
四个象限讨论一下即可. 
\end{jieda}


\item 若角$\alpha$是第三象限角, 则
\ding{172} $\sin\alpha+\cos\alpha<0$; 
\ding{173} $\tan\alpha-\sin\alpha>0$; 
\ding{174} $\cot\alpha\cdot\csc\alpha<0$; 
\ding{175} $\sin\alpha\cdot\sec\alpha>0$. 
其中正确的是\blkda{\ding{172}\ding{173}\ding{174}\ding{175}}.  

\item 已知角$\alpha$的终边过点$(3a-9,a+2)$, 且$\cos\alpha\leq 0$, $\sin\alpha>0$, 则$a$的取值范围是\blkda{$(-2,3]$}. 

\item 在$\triangle{ABC}$中, 有$\cos{A}\cdot\tan{B}\cdot\cot{C}<0$, 则这三角形是\blkda{钝角}三角形.
\begin{jieda}
首先没有直角, 否则要么积为零, 要么积不存在. \\
其次$\cos,\tan,\cot$在$(0,\frac{\pi}{2})$上都是正的, 在$(\frac{\pi}{2},\pi)$上都是负的, \\
故结合条件可得$\cos{A},\tan{B},\cot{C}$中一负两正, 于是钝角三角形. 
\end{jieda}

\item 下列各式为正号的是\dotfill(\kh{C})
\xx{$\cos2-\sin2$}
{$\cos2\cdot\sin2$}
{$\tan2\cdot\sec2$}
{$\sin2\cdot\tan2$}
\begin{jieda}
$2\in(\frac{\pi}{2},\pi)$, 故$\sin2>0$, $\cos2<0$, $\tan2<0$, $\cot2<0$, $\sec2<0$, $\csc2>0$. 
\end{jieda}

\item 设$y=\sin\alpha\cdot\cos\alpha\cdot\cot\alpha$, 且$\alpha$是象限角, 则$y$的符号为\dotfill(\kh{A})
\xx{恒正}{恒负}{可能为0}{不确定}

\item 已知实数$\alpha,\beta$满足$\abs{\cos\alpha-\cos\beta}=\abs{\cos\alpha}+\abs{\cos\beta}$, 且$\alpha\in(\frac{\pi}{2},\pi)$, 
则化简$\sqrt{(\cos\alpha-\cos\beta)^2}$的结果是\dotfill(\kh{C})
\xx{$\cos\alpha-\cos\beta$}
{$\abs{\cos\alpha}-\abs{\cos\beta}$}
{$\cos\beta-\cos\alpha$}
{$\abs{\cos\alpha}+\abs{\cos\beta}$}
\begin{jieda}
$\abs{\cos\alpha-\cos\beta}=\abs{\cos\alpha}+\abs{\cos\beta}\iff\cos\alpha\cos\beta\leq 0$. \\
又$\cos\alpha<0$, 故$\cos\beta\geq 0$. \\
故$E=\abs{\cos\alpha-\cos\beta}=\cos\beta-\cos\alpha$. 
\end{jieda}

\item 已知角$\alpha$的终边上一点$P$与点$A(-3,2)$关于$y$轴对称, 角$\beta$的终边上一点$Q$与点$A$关于原点对称, 求$2\sin\alpha+3\sin\beta$的值. 
\begin{jieda}
$A(-3,2)\xleftrightarrow{y\text{轴}}P(3,2)$, $A(-3,2)\xleftrightarrow{\text{原点}}Q(3,-2)$. \\
故$2\sin\alpha+3\sin\beta=2\cdot\dfrac{2}{\sqrt{13}}+3\cdot\dfrac{-2}{\sqrt{13}}=-\dfrac{2\sqrt{13}}{13}$. 
\end{jieda}


\item 已知角$\theta$终边上一点$P(-\sqrt{3},m)$, 且$\sin\theta=\dfrac{\sqrt{2}m}{4}$, 求$\cos\theta$和$\tan\theta$的值. 
\begin{jieda}
$\dfrac{m}{\sqrt{3+m^2}}=\dfrac{\sqrt{2}m}{4}$, 则$m=0,\pm\sqrt{5}$. 
\begin{itemize}
\item $m=0$时, $\cos\theta=-1$, $\tan\theta=0$; 
\item $m=\sqrt{5}$时, $\cos\theta=-\dfrac{\sqrt{6}}{4}$, $\tan\theta=-\dfrac{\sqrt{15}}{3}$; 
\item $m=-\sqrt{5}$时, $\cos\theta=-\dfrac{\sqrt{6}}{4}$, $\tan\theta=\dfrac{\sqrt{15}}{3}$. 
\end{itemize}
\end{jieda}

\end{enumerate}






\subsection{任意角的正弦、余弦、正切、余切(2)}
\zsd 特殊角三角比求值, 同角三角比关系(求值), 无三角函数线. 


\begin{tcolorbox}
\begin{itemize}
\item 定义: $\sin\alpha=\dfrac{y}{r}$, $\cos\alpha=\dfrac{x}{r}$. 
\item 令$r=1$, 则易得$\sin\alpha\in[-1,1]$, $\cos\alpha\in[-1,1]$, 且能得到$\sin\alpha,\cos\alpha$的变化规律. 
\end{itemize}
\end{tcolorbox}



\begin{enumerate}

\item 填表:

\begin{tabular}{|c|c|c|c|}\hline
$\alpha$ & $\sin\alpha$ & $\cos\alpha$ & $\tan\alpha$\\\hline
$\frac{2\pi}{3}$ & & & \\\hline
$\frac{5\pi}{4}$ & & & \\\hline
$\frac{5\pi}{6}$ & & & \\\hline
\end{tabular}


\item 填表:

\begin{tabular}{|c|c|c|c|c|c|}\hline
$\alpha$ & 点$P$的坐标 & ~~$\sin\alpha$~~ & ~~$\cos\alpha$~~ & ~~$\tan\alpha$~~ & ~~$\cot\alpha$~~ \\\hline
$\alpha=2k\pi(k\zz)$ & $(1,0)$ & 0 & 1 & 0 & 不存在\\\hline
& $(0,1)$ & & & & \\\hline
& $(-1,0)$ & & & & \\\hline
& $(0,-1)$ & & & & \\\hline
\end{tabular}


\item 计算: $\cos\dfrac{\pi}{3}-\tan\dfrac{\pi}{4}+\dfrac{3}{4}\tan^2\dfrac{7\pi}{6}-\sin\dfrac{5\pi}{6}+\cos^2\dfrac{\pi}{6}+\sin\dfrac{3\pi}{2}$=\blkda{$-1$}. 

\item 若$\sin\theta=\dfrac{m-3}{m+5}$, $\cos\theta=\dfrac{4-2m}{m+5}$都有意义, 则$m$=\blkda{0或8}. 
\begin{jieda}
$\sin^2\theta+\cos^2\theta=1$. 
\end{jieda}

\item 已知$\sin\alpha=\dfrac{4}{5}$, $\alpha\in(0,\pi)$, 则$\tan\alpha$等于\blkda{$\pm\dfrac{4}{3}$}. 
\begin{jieda}
三角比的值分为两部分: 符号与绝对值. 结合定义, 直接赋值: $y=4,r=5$, 于是$x=\pm3$. 
\end{jieda}

\item 若$\cot(\sin\theta)\cdot\tan(\cos\theta)>0$, 则$\theta$是第\blkda{I、III}象限角. 
\begin{jieda}
四个象限讨论一下即可. $\cos\theta,\sin\theta\in[-1,1]\subseteq(-\frac{\pi}{2},\frac{\pi}{2})$. 
\end{jieda}

\item 若$\tan\alpha=\sqrt3$, $\pi<\alpha<\dfrac{3\pi}{2}$, 那么$\cos\alpha-\sin\alpha$的值是\dotfill(\kh{B})
\xx{$-\dfrac{1+\sqrt3}{2}$}
{$\dfrac{-1+\sqrt3}{2}$}
{$\dfrac{1-\sqrt3}{2}$}
{$\dfrac{1+\sqrt3}{2}$}
\begin{jieda}
由条件可取终边上一点$(-1,-\sqrt{3})$. 
\end{jieda}

\item 已知$\sin\alpha=m(\abs{m}<1)$, $\dfrac{\pi}{2}<\alpha<\dfrac{3\pi}{2}$, 那么$\tan\alpha=$\dotfill(\kh{B})
\xx{$\dfrac{m}{\sqrt{1-m^2}}$}
{$-\dfrac{m}{\sqrt{1-m^2}}$}
{$\pm\dfrac{m}{\sqrt{1-m^2}}$}
{$\pm\dfrac{\sqrt{1-m^2}}{m}$}
\begin{jieda}
因为$\dfrac{\pi}{2}<\alpha<\dfrac{3\pi}{2}$, 故终边上点的横坐标为负. \\
$m=\dfrac{m}{1}$, 故终边上取点$(-\sqrt{1-m^2},m)$, 于是立得. 
\end{jieda}

\item 已知$\cot\alpha=\dfrac{1}{3}$, 求$\sin\alpha$、$\cos\alpha$及$\tan\alpha$. 
\begin{jieda}
$\alpha$为第一象限角时, $\sin\alpha=\dfrac{3}{\sqrt{10}}$, $\cos\alpha=\dfrac{1}{\sqrt{10}}$, $\tan\alpha=3$; \\
$\alpha$为第三象限角时, $\sin\alpha=-\dfrac{3}{\sqrt{10}}$, $\cos\alpha=-\dfrac{1}{\sqrt{10}}$, $\tan\alpha=3$. 
\end{jieda}

\item 已知$\sin\alpha=\dfrac{m^2-1}{m^2+1}(m>0)$, 求$\cos\alpha$与$\tan\alpha$的值. 
\begin{jieda}
$\alpha$为第I象限角时, $\cos\alpha=\dfrac{2m}{m^2+1}$, $\tan\alpha=\dfrac{m^2-1}{2m}$; \\
$\alpha$为第II象限角时, $\cos\alpha=-\dfrac{2m}{m^2+1}$, $\tan\alpha=-\dfrac{m^2-1}{2m}$. 
\end{jieda}
\end{enumerate}







\subsection{任意角的正弦、余弦、正切、余切(3)}
\zsd 同角三角比关系: 求值和证明

\begin{enumerate}
\item 若$\theta$是锐角, $\sin\theta-\cos\theta=\dfrac{1}{2}$, 则$\sin^3\theta-\cos^3\theta$=\blkda{$\dfrac{11}{16}$}. 
\begin{jieda}
围绕$(\sin\theta\pm\cos\theta)^2=1\pm2\sin\theta\cos\theta$展开. 
\begin{align*}
E & =(\sin\theta-\cos\theta)(\sin^2\theta+\sin\theta\cos\theta+\cos^2\theta)\\
& = (\sin\theta-\cos\theta)\left(1+\dfrac{1-(\sin\theta-\cos\theta)^2}{2}\right)\\
& = \dfrac{1}{2}\cdot\dfrac{3-\frac{1}{4}}{2}\\
& = \dfrac{11}{16}.
\end{align*}
\end{jieda}

\item 设$\cot{x}=2$, 则$\dfrac{2\cos{x}-4\sin{x}}{5\cos{x}+3\sin{x}}$=\blkda{0}, 
$3\sin^2x-4\cos^2x$=\blkda{$-\dfrac{13}{5}$}. 
\begin{jieda}
$\dfrac{2\cos{x}-4\sin{x}}{5\cos{x}+3\sin{x}}=\dfrac{2\cot{x}-4}{5\cot{x}+3}=0$. \\
$3\sin^2x-4\cos^2x=\dfrac{3\sin^2x-4\cos^2x}{\cos^2x+\sin^2x}=\dfrac{3-4\cot^2x}{\cot^2x+1}=-\dfrac{13}{5}$. 
\end{jieda}

\item 已知$\tan\alpha=\dfrac{2ab}{a^2-b^2}$, 其中$a>b>0$, $\alpha\in(0,\frac{\pi}{2})$, 则$\sin\alpha$=\blkda{$\dfrac{2ab}{a^2+b^2}$}.
\begin{jieda}
$\sin\alpha=\dfrac{2ab}{\sqrt{(2ab)^2+(a^2-b^2)^2}}$. 
\end{jieda}

\item 已知$\dfrac{\cos\alpha}{1+\sin\alpha}=-\dfrac{1}{2}$, 则$\dfrac{\cos\alpha}{\sin\alpha-1}$=\blkda{2}. 
\begin{jieda}
结合$\cos^2\alpha+\sin^2\alpha=1$解得$(\cos\alpha,\sin\alpha)=(-\frac{4}{5},\frac{3}{5})$(另一组增根已舍去), 
代入所求即得. 
\end{jieda}

\item 用列举法写出集合$A=\left\{y\mid y=\dfrac{1}{\cos\alpha\sqrt{1+\tan^2\alpha}}
+\dfrac{2\tan\alpha}{\sqrt{\sec^2\alpha-1}}\right\}$=\blkda{$\{\pm1,\pm3\}$}.
\begin{jieda}
$y=\dfrac{\abs{\cos\alpha}}{\cos\alpha}+\dfrac{2\tan\alpha}{\abs{\tan\alpha}}$. 
\end{jieda}

\item 使函数$y=\sqrt{\sin{x}}+\sqrt{16-x^2}$有意义的$x$的取值范围是\blkda[3]{$[-4,-\pi]\cup[0,\pi]$}; 
使函数$y=\lg\sin(\cos{x})$有意义的$x$的取值范围是\blkda[5]{$\left(-\dfrac{\pi}{2}+2k\pi,\dfrac{\pi}{2}+2k\pi\right)(k\zz)$}. 
\begin{jieda}
$\sin{x}\geq 0\land 16-x^2\geq 0$, 故$x\in[-4,-\pi]\cup[0,\pi]$. \\
$\sin(\cos{x})>0\land\cos{x}\in[-1,1]\iff\cos{x}\in(0,1]\iff x\in\left(-\dfrac{\pi}{2}+2k\pi,\dfrac{\pi}{2}+2k\pi\right)(k\zz)$. 
\end{jieda}

\item 若角$\alpha$的终边落在直线$x+y=0$上, 则$\dfrac{\sin\alpha}{\sqrt{1-\sin^2\alpha}}+\dfrac{\sqrt{1-\cos^2\alpha}}{\cos\alpha}$的值等于\dotfill(\kh{D})
\xx{2}{$-2$}{$-2$或2}{0}
\begin{jieda}
考虑一下符号基本上就可以把答案选出来了. 
\end{jieda}

\item 已知$\theta\in\left(0,\dfrac{\pi}{2}\right)$, 且$\log_{\tan\theta+\cot\theta}\sin\theta=-\dfrac{3}{4}$, 
则$\log_{\tan\theta}\cos\theta=$\dotfill(\kh{C})
\xx{2}{$-2$}{$\dfrac{1}{2}$}{$-\dfrac{1}{2}$}
\begin{jieda}
$(\tan\theta+\cot\theta)^{-\frac{3}{4}}=\sin\theta$, 
即$1=\sin^4\theta(\tan\theta+\cot\theta)^3=\sin^4\theta\cdot\left(\dfrac{\sin\theta}{\cos\theta}+\dfrac{\cos\theta}{\sin\theta}\right)^3$. \\
于是$1=\dfrac{\sin^4\theta}{\sin^3\theta\cos^3\theta}=\dfrac{\sin\theta}{\cos^3\theta}=\dfrac{\tan\theta}{\cos^2\theta}$. \\
故$\cos^2\theta=\tan\theta$, 所以$\log_{\tan\theta}\cos\theta=\dfrac{1}{2}$. 
\end{jieda}


\item 已知$\sin\alpha+\cos\alpha=\dfrac{1}{5}(0<\alpha<\pi)$, 求: 
\begin{enumerate}[(1)]
\item $\sin\alpha-\cos\alpha$; 
\item $\sin^3\alpha+\cos^3\alpha$; 
\item $\sin^4\alpha+\cos^4\alpha$. 
\end{enumerate}
\begin{jieda}
结合$\sin^2\alpha+\cos^2\alpha=1$解得$(\cos\alpha,\sin\alpha)=(-\frac{3}{5},\frac{4}{5})$. (另一组解已舍去)\\
\begin{enumerate*}[(1)]
\item $\dfrac{7}{5}$; 
\item $\dfrac{37}{125}$; 
\item $\dfrac{337}{625}$. 
\end{enumerate*}
\end{jieda}

\item 已知$\tan\alpha=\sqrt{2}$, 求下列各式的值:
\begin{enumerate}[(1)]
\item $\sin\alpha+2\cos\alpha$; 
\item $\dfrac{\cos\alpha-5\sin\alpha}{3\cos\alpha+\sin\alpha}$; 
\item $\dfrac{\sin^2\alpha-\sin\alpha\cos\alpha-3\cos^2\alpha}{5\sin\alpha\cos\alpha+\sin^2\alpha+1}$; 
\item $2\sin^2\alpha-\sin\alpha\cos\alpha+\cos^2\alpha$. 
\end{enumerate}
\begin{jieda}
齐次化. \\
\begin{enumerate*}[(1)]
\item $\pm\dfrac{\sqrt{6}+2\sqrt{3}}{3}$; 
\item $\dfrac{13-16\sqrt2}{7}$; 
\item $-\dfrac{1}{5}$; 
\item $\dfrac{5-\sqrt{2}}{3}$. 
\end{enumerate*}
\end{jieda}

\item 证明下列恒等式: 
\begin{enumerate}[(1)]
\item $\dfrac{1+2\sin{x}\cos{x}}{\cos^2x-\sin^2x}=\dfrac{1+\tan{x}}{1-\tan{x}}$; 
\item $\dfrac{1-\sin^6x-\cos^6x}{1-\sin^4x-\cos^4x}=\dfrac{3}{2}$;
\item $\dfrac{\tan\alpha\cdot\sin\alpha}{\tan\alpha-\sin\alpha}=\dfrac{\tan\alpha+\sin\alpha}{\tan\alpha\cdot\sin\alpha}$; 
\item $\dfrac{\cos\alpha}{1+\sin\alpha}-\dfrac{\sin\alpha}{1+\cos\alpha}=\dfrac{2(\cos\alpha-\sin\alpha)}{1+\sin\alpha+\cos\alpha}$. 
\end{enumerate}
\begin{jieda}
\begin{enumerate}[(1)]
\item $\LHS=\dfrac{(\sin{x}+\cos{x})^2}{(\cos{x}-\sin{x})(\cos{x}+\sin{x})}=\dfrac{\cos{x}+\sin{x}}{\cos{x}-\sin{x}}$. 

\item ``$\sin^2\theta+\cos^2\theta=1$''. 
\begin{align*}
\LHS & = \dfrac{1-(\sin^2x+\cos^2x)(\sin^4x-\sin^2x\cos^2x+\cos^4x)}{1-(\sin^2x+\cos^2x)^2+2\sin^2x\cos^2x}\\
& = \dfrac{1-(\sin^2x+\cos^2x)^2+3\sin^2x\cos^2x}{2\sin^2x\cos^2x}\\
& = \dfrac{3\sin^2x\cos^2x}{2\sin^2x\cos^2x}
\end{align*}

\item $\iff\tan^2\alpha\sin^2\alpha=\tan^2\alpha-\sin^2\alpha\iff\sin^2\alpha=\tan^2\alpha-\tan^2\alpha\sin^2\alpha$. 

\item 通分, 配凑. 
\begin{align*}
\LHS & = \dfrac{\cos\alpha(1+\cos\alpha)-\sin\alpha(1+\sin\alpha)}{(1+\sin\alpha)(1+\cos\alpha)}\\
& = \dfrac{(\cos\alpha-\sin\alpha)(1+\cos\alpha+\sin\alpha)}{(1+\sin\alpha)(1+\cos\alpha)}\\
& = \dfrac{2(\cos\alpha-\sin\alpha)}{1+\sin\alpha+\cos\alpha}\cdot\dfrac{(1+\cos\alpha+\sin\alpha)^2}{2(1+\sin\alpha)(1+\cos\alpha)}\\
& = \RHS\cdot\dfrac{1+\cos^2\alpha+\sin^2\alpha+2\cos\alpha+2\sin\alpha+2\cos\alpha\sin\alpha}{2(1+\cos\alpha+\sin\alpha+\cos\alpha\sin\alpha)}\\
& = \RHS
\end{align*}
\end{enumerate}
\end{jieda}


\item
\begin{enumerate}
\item 已知$3\sin^2\alpha+2\sin^2\beta=5\sin\alpha$, 求$\sin^2\alpha+\sin^2\beta$的范围; 
\item 已知$6\sin3\alpha-\cos^22\beta=6$, 求角$\alpha$的值. 
\end{enumerate}
\begin{jieda}
\begin{enumerate}
\item $2\sin^2\beta=5\sin\alpha-3\sin^2\alpha\in[0,2]\land\sin\alpha\in[-1,1]$, 则$\sin\alpha\in[0,\frac{2}{3}]\cup\{1\}$. \\
故$E=\sin^2\alpha+\frac{5}{2}\sin\alpha-\frac{3}{2}\sin^2\alpha=-\frac{1}{2}\sin^2\alpha+\frac{5}{2}\sin\alpha\in[0,\frac{13}{9}]\cup\{2\}$. 

\item $6\sin3\alpha-6=\cos^22\beta$. \\
$\LHS\in[-12,0],\RHS\in[0,1]$, 故$6\sin3\alpha-6=\cos^22\beta=0$, \\
于是$3\alpha=2k\pi+\dfrac{\pi}{2}$, 即$\alpha=\dfrac{2k\pi}{3}+\dfrac{\pi}{6}(k\zz)$. 
\end{enumerate}
\end{jieda}
\end{enumerate}




\subsection{诱导公式(1)}
\zsd ``符号看象限, 奇变偶不变''之``偶''.

\begin{enumerate}
\item 若$\cos(\pi-\alpha)=\dfrac{\sqrt5}{3}$, 且$\alpha\in\left(\dfrac{\pi}{2},\pi\right)$, 则$\sin(\pi+\alpha)$=\blkda{$-\dfrac{2}{3}$}. 

\item 已知$\sin(\alpha-\pi)=\dfrac{2}{3}$, 且$\alpha\in\left(-\dfrac{\pi}{2},0\right)$, 则$\tan\alpha$=\blkda{$-\dfrac{2}{\sqrt{5}}$}. 

\item 已知$\sin\left(\alpha-\dfrac{2\pi}{3}\right)=\dfrac{1}{4}$, 则$\sin\left(\alpha+\dfrac{\pi}{3}\right)$=\blkda{$-\dfrac{1}{4}$}. 


\item 已知$\cos\left(\dfrac{\pi}{6}+\theta\right)=\dfrac{\sqrt{3}}{3}$, 则$\cos\left(\dfrac{5\pi}{6}-\theta\right)$=\blkda{$\dfrac{-\sqrt{3}}{3}$}. 

\item 已知$\cos(11\pi-3)=p$, 用$p$表示$\tan(-3)$=\blkda{$\dfrac{\sqrt{1-p^2}}{p}$}.  

\item 若$\tan\alpha=-2$, 则$\sin(\alpha-\pi)\cdot\cos(\pi+\alpha)$=\blkda{$-\dfrac{2}{5}$}. 

\item $\sin^2(\alpha+\pi)-\cos(\alpha+\pi)\cos(-\alpha)+1$的值是\dotfill(\kh{D})
\xx{1}{$2\sin^2\alpha$}{0}{2}

\item 在$\triangle{ABC}$中, 给出下列四个式子: 
\ding{172} $\sin(A+B)+\sin{C}$, 
\ding{173} $\cos(A+B)+\cos{C}$, 
\ding{174} $\sin(2A+2B)+\sin2C$, 
\ding{175} $\cos(2A+2B)+\cos2C$, 其中恒为常数的是\dotfill(\kh{B})
\xx{\ding{172}\ding{173}}{\ding{173}\ding{174}}{\ding{174}\ding{175}}{以上都不对}


\item 函数式$\sqrt{1+2\sin(\pi-2)\cos(\pi+2)}$化简的结果是\dotfill(\kh{A})
\xx{$\sin2-\cos2$}
{$\pm(\sin2-\cos2)$}
{$\cos2-\sin2$}
{以上结论都不对}

\item 当$n\zz$时, 在
\ding{172} $\sin(n\pi+\frac{\pi}{3})$; 
\ding{173} $\sin(2n\pi\pm\frac{\pi}{3})$; 
\ding{174} $\sin(n\pi+(-1)^n\frac{\pi}{3})$; 
\ding{175} $\cos(2n\pi\pm\frac{\pi}{6})$; 
\ding{176} $\cos(n\pi+(-1)^n\frac{\pi}{6})$中与$\sin\frac{\pi}{3}$相等的是\dotfill(\kh{A})
\xx{\ding{174}\ding{175}}
{\ding{172}\ding{174}}
{\ding{173}\ding{175}\ding{176}}
{\ding{174}\ding{175}\ding{176}}


\item 已知$\sin(-\pi+\alpha)+2\cos(3\pi-\alpha)=0$, 计算: 
\begin{enumerate}[(1)]
\item $\dfrac{2\sin\alpha-\cos\alpha}{\sin\alpha+3\cos\alpha}$; 
\item $\sin^2\alpha+\sin\alpha\cos\alpha-3\cos^2\alpha$. 
\end{enumerate}
\begin{jieda}
条件可化为$-\sin\alpha-2\cos\alpha=0$, 故$\tan\alpha=-2$. 
\begin{enumerate}[(1)]
\item $E=\dfrac{2\tan\alpha-1}{\tan\alpha+3}=-5$. 
\item $E=\dfrac{\sin^2\alpha+\sin\alpha\cos\alpha-3\cos^2\alpha}{\sin^2\alpha+\cos^2\alpha}=\dfrac{\tan^2\alpha+\tan\alpha-3}{\tan^2\alpha+1}=-\dfrac{1}{5}$
\end{enumerate}
\end{jieda}

\item 设$f(\theta)=\dfrac{2\cos^3\theta+\sin^2(2\pi-\theta)+\cos(-\theta)-3}{2+2\cos^2(\pi+\theta)+\cos(2\pi-\theta)}$, 求$f(\frac{\pi}{3})$的值. 
\begin{jieda}
$f(\theta)=\dfrac{2\cos^3\theta+1-\cos^2\theta+\cos\theta-3}{2\cos^2\theta+\cos\theta-2}=\cos\theta-1$. \\
故$f(\frac{\pi}{3})=-\dfrac{1}{2}$. 
\end{jieda}

\item 已知$\sin(3\pi-\alpha)=\sqrt2\sin(2\pi+\beta)$, $\sqrt3\cos(-\alpha)=-\sqrt2\cos(\pi+\beta)$, 
且$0<\alpha<\pi$, $0<\beta<\pi$, 求$\sin\alpha,\sin\beta$. 
\begin{jieda}
$\begin{cases}
\sin\alpha=\sqrt{2}\sin\beta\\
\sqrt{3}\cos\alpha=\sqrt{2}\cos\beta
\end{cases}$, 两式两边平方再相加可得$\sin^2\alpha+3\cos^2\alpha=2$, 故$\sin^2\alpha=\dfrac{1}{2}$. \\
又$\alpha\in(0,\pi)$, 故$\sin\alpha=\dfrac{\sqrt{2}}{2}$. 
代入第一式可得$\sin\beta=\dfrac{1}{2}$. 
% $\sin\alpha=\dfrac{\sqrt{2}}{2}$, $\sin\beta=\dfrac{1}{2}$. 
\end{jieda}

\item 已知$f(x)=\dfrac{\sin(n\pi-x)\cos(n\pi-x)}{\cos((n+1)\pi-x)}\cdot\tan(x-n\pi)\cot(n\pi-x)(n\zz)$, 求$f(\frac{7\pi}{6})$. 
\begin{jieda}
$n$为偶数时, $f(x)=-\sin{x}$; $n$为奇数时, $f(x)=\sin{x}$. \\
$f(\frac{7\pi}{6})=\begin{cases}
-\frac{1}{2}, & n\text{为奇数}\\
\frac{1}{2}, & n\text{为偶数}
\end{cases}$. 
\end{jieda}

\end{enumerate}




\subsection{诱导公式(2)}
\zsd ``符号看象限, 奇变偶不变''之``奇''.

\begin{enumerate}
\item 已知$\sin(\frac{\pi}{2}-\alpha)=\frac{3}{5}$, 则$\cos(\pi+\alpha)$的值为\blkda{$-\dfrac{3}{5}$}. 

\item 已知$\cos(\frac{\pi}{2}+\varphi)=\dfrac{\sqrt3}{2}$, 且$\abs{\varphi}<\dfrac{\pi}{2}$, 则$\tan\varphi$=\blkda{$-\sqrt{3}$}. 
\begin{jieda}
$-\sin\varphi=\dfrac{\sqrt{3}}{2}$, 故$\sin\varphi=-\dfrac{\sqrt{3}}{2}$. \\
又$\abs{\varphi}<\dfrac{\pi}{2}$, 故$-\dfrac{\pi}{2}<\varphi<0$, 于是$\tan\varphi=-\sqrt{3}$. 
\end{jieda}

\item 如果$\cos(\pi+\alpha)=-\dfrac{1}{3}$, 那么$\sin\left(\frac{3\pi}{2}-\alpha\right)$=\blkda{$-\dfrac{1}{3}$}. 
\begin{jieda}
$-\dfrac{1}{3}=\cos(\pi+\alpha)=-\cos\alpha$, 故$\cos\alpha=\dfrac{1}{3}$. \\
于是$E=-\cos\alpha=-\dfrac{1}{3}$. 
\end{jieda}

\item 已知函数$f(x)=\sqrt2\cos(x-\frac{\pi}{12})$, 若$\cos\theta=\frac{3}{5}$, $\theta\in(\frac{3\pi}{2},2\pi)$, 
则$f(\theta-\frac{5\pi}{12})$=\blkda{$-\frac{4\sqrt{2}}{5}$}. 
\begin{jieda}
$f(\theta-\frac{5\pi}{12})=\sqrt2\cos(\theta-\frac{5\pi}{12}-\frac{\pi}{12})=\sqrt{2}\sin\theta=-\frac{4\sqrt{2}}{5}$
\end{jieda}


\item 已知$\sin\beta=\dfrac{1}{3}$, $\sin(\alpha+\beta)=1$, 则$\sin(2\alpha+\beta)$=\blkda{$\dfrac{1}{3}$}. 
\begin{jieda}
$\sin(\alpha+\beta)=1\iff\alpha+\beta=2k\pi+\dfrac{\pi}{2}\iff\alpha=\dfrac{\pi}{2}-\beta+2k\pi(k\zz)$. \\
故$\sin(2\alpha+\beta)=\sin(\pi-2\beta+4k\pi+\beta)=\sin\beta=\dfrac{1}{3}$. 
\end{jieda}


\item 设$f(n)=\sin\left(\frac{n\pi}{4}+\alpha\right)$, 则$f(n)\cdot f(n+4)+f(n+2)\cdot f(n+6)$=\blkda{$-1$}. 
\begin{jieda}
令$\beta=\dfrac{n\pi}{4}+\alpha$, 故$E=\sin\beta\sin(\pi+\beta)+\sin(\frac{\pi}{2}+\beta)\sin(\pi+\frac{\pi}{2}+\beta)=-\sin^2\beta-\cos^2\beta=-1$. 
\end{jieda}

\item 已知$\sin(\frac{\pi}{2}+\theta)+3\cos(\pi-\theta)=\sin(-\theta)$, 则$\sin\theta\cos\theta+\cos^2\theta=$\dotfill(\kh{C})
\xx{$\dfrac{1}{5}$}
{$\dfrac{2}{5}$}
{$\dfrac{3}{5}$}
{$\dfrac{\sqrt5}{5}$}
\begin{jieda}
$\cos\theta-3\cos\theta=-\sin\theta$, 故$\tan\theta=2$. \\
$E=\dfrac{\sin\theta\cos\theta+\cos^2\theta}{\sin^2\theta+\cos^2\theta}=\dfrac{\tan\theta+1}{\tan^2\theta+1}=\dfrac{3}{5}$. 
\end{jieda}

\item 已知点$P(\sin(\pi+\theta),\sin(\frac{3\pi}{2}-\theta))$在第三象限, 则角$\theta$所在的象限是\dotfill(\kh{A})
\xx{第一象限}{第二象限}{第三象限}{第四象限}
\begin{jieda}
$\begin{cases}
\sin(\pi+\theta)<0\\
\sin(\frac{3\pi}{2}-\theta)<0
\end{cases}$, 
于是$\begin{cases}
\sin\theta>0\\
\cos\theta>0
\end{cases}$. 
\end{jieda}


\item 已知角$\alpha$的终边上一点的坐标为$(\sin\frac{4\pi}{3},\cos\frac{4\pi}{3})$, 则角$\alpha$的最小正值为\dotfill(\kh{A})
\xx{$\dfrac{7\pi}{6}$}
{$\dfrac{11\pi}{6}$}
{$\dfrac{5\pi}{6}$}
{$\dfrac{4\pi}{3}$}
\begin{jieda}
法一: 点$(-\frac{\sqrt{3}}{2},-\frac{1}{2})$, 故所求为$\dfrac{7\pi}{6}$. \\
法二: $f(\frac{\pi}{2}-\theta)=g(\theta)$. 
\end{jieda}

\item 已知角$\alpha$的顶点与原点$O$重合, 始边与$x$轴的正半轴重合, 将$\alpha$的终边按顺时针方向旋转$\frac{\pi}{2}$后, 
过点$P(\frac{3}{5},\frac{4}{5})$, 则$\cos\alpha$等于\dotfill(\kh{A})
\xx{$-\dfrac{4}{5}$}
{$\dfrac{4}{5}$}
{$-\dfrac{3}{5}$}
{$\dfrac{3}{5}$}
\begin{jieda}
$\beta=\alpha-\frac{\pi}{2}$, 则$\alpha=\beta+\frac{\pi}{2}$. 
故$\cos\alpha=\cos(\beta+\frac{\pi}{2})=-\sin\beta=-\frac{4}{5}$. 
\end{jieda}

\item 已知$f(x)=\dfrac{\cos(\frac{\pi}{2}+x)\cos(-x)\sin(\frac{3\pi}{2}-x)}{\sin(-\pi-x)\cos(2\pi-x)}$. 
\begin{enumerate}[(1)]
\item 化简$f(x)$; 
\item 若$x$是第三象限角, 且$\tan{x}=2$, 求$f(x)$的值. 
\end{enumerate}
\begin{jieda}
\begin{enumerate}[(1)]
\item $f(x)=\dfrac{-\sin{x}\cos{x}(-\cos{x})}{\sin{x}\cos{x}}=\cos{x}$. 
\item $f(x)=-\dfrac{1}{\sqrt{5}}$. 
\end{enumerate}
\end{jieda}


\item 已知$\sin\alpha$是方程$3x^2-10x-8=0$的根, 且$\alpha$为第三象限角, \\
求$\dfrac{\sin(\alpha+\frac{3\pi}{2})\sin(\frac{3\pi}{2}-\alpha)\tan^2(2\pi-\alpha)\tan(\pi-\alpha)}
{\cos(\frac{\pi}{2}-\alpha)\cos(\frac{\pi}{2}+\alpha)}$的值. 
\begin{jieda}
易得$\sin\alpha=-\dfrac{2}{3}$. 
故
\begin{align*}
E & = \dfrac{-\cos\alpha(-\cos\alpha)\tan^2\alpha(-\tan\alpha)}{\sin\alpha(-\sin\alpha)}\\
& = \tan\alpha\\
& = \dfrac{2}{\sqrt{5}}.
\end{align*}
\end{jieda}

\end{enumerate}




\subsection{已知正弦、余弦、正切求角}

\begin{tcolorbox}
\zsd 
\begin{itemize}
\item 若$\sin{x}=\sin\alpha$, 则$x=2k\pi+\alpha$或$x=2k\pi+\pi-\alpha$, 即$x=k\pi+(-1)^k\alpha, k\zz$. 
\item 若$\cos{x}=\cos\alpha$, 则$x=2k\pi\pm\alpha, k\zz$. 
\item 若$\tan{x}=\tan\alpha$, 则$x=k\pi+\alpha, k\zz$. 
\end{itemize}
\end{tcolorbox}

\begin{enumerate}
\item 已知$\sin2x=-\dfrac{\sqrt2}{2}, x\in[-\pi,\pi]$, 则$x$=\blkda[3]{$-\dfrac{3\pi}{8},-\dfrac{\pi}{8},\dfrac{5\pi}{8},\dfrac{7\pi}{8}$}. 
\begin{jieda}
$2x\in[-2\pi,2\pi]$, 故$2x=-\dfrac{3\pi}{4},-\dfrac{\pi}{4},\dfrac{5\pi}{4},\dfrac{7\pi}{4}$, 
故$x=-\dfrac{3\pi}{8},-\dfrac{\pi}{8},\dfrac{5\pi}{8},\dfrac{7\pi}{8}$. 
\end{jieda}

\item 已知$\tan\left(2x+\dfrac{\pi}{6}\right)=-1$, 则$x$=\blkda[3]{$\dfrac{k\pi}{2}-\dfrac{5\pi}{24}(k\zz)$}. 
\begin{jieda}
$2x+\dfrac{\pi}{6}=k\pi-\dfrac{\pi}{4}$, 故$x=\dfrac{k\pi}{2}-\dfrac{5\pi}{24}(k\zz)$
\end{jieda}

\item 方程$2\sin^2x+3\cos{x}=0$, $x\in[-2\pi,2\pi]$的解集为\blkda[3]{$\left\{\pm\dfrac{4\pi}{3},\pm\dfrac{2\pi}{3}\right\}$} .
\begin{jieda}
$\iff 2-2\cos^2x+3\cos{x}=0$, 故$\cos{x}=-\dfrac{1}{2}$(2已舍去), 于是$x=2k\pi\pm\dfrac{2\pi}{3}(k\zz)$. \\
又$x\in[-2\pi,2\pi]$, 故$x=-\dfrac{4\pi}{3},-\dfrac{2\pi}{3},\dfrac{2\pi}{3},\dfrac{4\pi}{3}$. 
\end{jieda}

\item 已知$\sin\dfrac{\alpha}{2}=-\dfrac{\sqrt{3}}{2}$, 且$\alpha$是第二象限角, 则角$\alpha$=\blkda[3]{$4k\pi-\dfrac{4\pi}{3}(k\zz)$}. 
\begin{jieda}
$\dfrac{\alpha}{2}=2k\pi-\dfrac{\pi}{3}$或$2k\pi-\dfrac{2\pi}{3}$, 故$\alpha=4k\pi-\dfrac{2\pi}{3}$或$4k\pi-\dfrac{4\pi}{3}$. \\
又$\alpha$是第二象限角, 故$\alpha=4k\pi-\dfrac{4\pi}{3}(k\zz)$
\end{jieda}


\item 方程$\sin\left(x+\dfrac{\pi}{6}\right)-\cos{x}=0$的解集为\blkda[4]{$\left\{x\mid x=k\pi+\dfrac{\pi}{6},k\zz\right\}$}. 
\begin{jieda}
$\sin\alpha=\cos\beta\iff \sin\alpha=\sin\left(\dfrac{\pi}{2}-\beta\right)\iff
\alpha=\dfrac{\pi}{2}-\beta+2k\pi$或$\alpha=\dfrac{\pi}{2}+\beta+2k\pi$. 
\begin{itemize}
\item $x+\dfrac{\pi}{6}=\dfrac{\pi}{2}-x+2k\pi\iff x=k\pi+\dfrac{\pi}{6}(k\zz)$; 
\item $x+\dfrac{\pi}{6}=\dfrac{\pi}{2}+x+2k\pi\iff 2k\pi+\dfrac{\pi}{3}=0$, 无解. 
\end{itemize}
\end{jieda}


\item 已知集合$A=\left\{\alpha\mid 2\sin\alpha-1\geq 0\right\}$, $B=\left\{\alpha\mid\sqrt2\cos\alpha+1\geq 0\right\}$, \\
则$A\cap B$=\blkda[7]{$\left\{\alpha\mid 2k\pi+\dfrac{\pi}{6}\leq\alpha\leq 2k\pi+\dfrac{3\pi}{4},k\zz\right\}$}. 
\begin{jieda}
$\sin\alpha\geq\dfrac{1}{2}\iff 2k\pi+\dfrac{\pi}{6}\leq\alpha\leq 2k\pi+\dfrac{5\pi}{6}$, \\
$\cos\alpha\geq-\dfrac{1}{\sqrt{2}}\iff 2k\pi-\dfrac{3\pi}{4}\leq\alpha\leq 2k\pi+\dfrac{3\pi}{4}$. \\
于是$A\cap B=\left\{\alpha\mid 2k\pi+\dfrac{\pi}{6}\leq\alpha\leq 2k\pi+\dfrac{3\pi}{4},k\zz\right\}$. 
\end{jieda}


\item ``$\alpha=\dfrac{\pi}{6}+2k\pi(k\zz)$''是``$\cos2\alpha=\dfrac{1}{2}$''的\dotfill(\kh{A})
\xx{充分而不必要条件}{必要而不充分条件}{充分必要条件}{既不充分也不必要条件}


\item 若三角方程$\sin{x}=0$与$\sin2x=0$的解集分别为$E,F$, 则\dotfill(\kh{A})
\xx{$E\subsetneqq F$}{$E\supsetneqq F$}{$E=F$}{$E\cap F=\varnothing$}


\item 方程$\sqrt{1-\cos^2x}=\cos{x}$的解集是\dotfill(\kh{D})
\xx
{$\left\{x\mid x=k\pi\pm\dfrac{\pi}{4}, k\zz\right\}$}
{$\left\{x\mid x=k\pi+\dfrac{\pi}{4}, k\zz\right\}$} 
{$\left\{x\mid x=k\pi-\dfrac{\pi}{4}, k\zz\right\}$}
{$\left\{x\mid x=2k\pi\pm\dfrac{\pi}{4}, k\zz\right\}$}


\item 已知$\theta\in[0,2\pi]$, $\sin\theta$和$\cos\theta$是方程$x^2-kx+k+1=0$的两个根, 求$\theta$. 
\begin{jieda}
$\Delta\geq 0$, 则$k\in(-\infty,2-2\sqrt{2}]\cup[2+2\sqrt{2},+\infty)$. \\
于是$\begin{cases}
\sin\theta+\cos\theta=k\\
\sin\theta\cos\theta=k+1
\end{cases}$, 
而$(\sin\theta+\cos\theta)^2=1+2\sin\theta\cos\theta$, \\
故$k^2=1+2(k+1)$, 解得$k=-1$或3. \\
当$k=3$时, $\sin\theta\cos\theta=4>1$, 不符合条件, 故舍去; \\
当$k=-1$时, 解得$(\cos\theta,\sin\theta)=(0,-1)$或$(-1,0)$. 又$\theta\in[0,2\pi]$, 故$\theta=\dfrac{3\pi}{2},\pi$. 
\end{jieda}


\item 已知集合$A=\left\{-\dfrac{1}{2},3,\cos{x}\right\}$, $B=\left\{\dfrac{\sqrt{3}}{2},\sin{x}\right\}$. 
\begin{enumerate}[(1)]
\item 当$B\subsetneqq A$时, 求$x$的值;
\item 当$A\cap B=\{y\}$时, 求$x$和$y$的值. 
\end{enumerate}
\begin{jieda}
\begin{enumerate}[(1)]
\item 由题意可得$\begin{cases}
-\frac{1}{2}=\sin{x}\\
\cos{x}=\frac{\sqrt{3}}{2}
\end{cases}$, 故$x=2k\pi+\dfrac{11\pi}{6}(k\zz)$. 

\item 若$y=\dfrac{\sqrt{3}}{2}$, 则$\cos{x}=\dfrac{\sqrt{3}}{2}\land\sin{x}\neq-\dfrac{1}{2}$, 故$x=2k\pi+\dfrac{\pi}{6}(k\zz)$. \\
若$y=\sin{x}$, 则$y=\cos{x}$或$y=-\dfrac{1}{2}$, 即 (i)$x=k\pi+\dfrac{\pi}{4}$, 或(ii)$x=2k\pi+\dfrac{7\pi}{6}$或$x=2k\pi+\dfrac{11\pi}{6}$. \\
检验可得$x=2k\pi+\dfrac{11\pi}{6}$不符合条件. \\
综上, $(x,y)=(2k\pi+\frac{\pi}{6},\frac{\sqrt{3}}{2})$, $(2k\pi+\frac{\pi}{4},\frac{\sqrt{2}}{2})$, 
$(2k\pi+\frac{5\pi}{4},-\frac{\sqrt{2}}{2})$, $(2k\pi+\frac{7\pi}{6},-\frac{1}{2})$, 其中$k\zz$. 
\end{enumerate}
\end{jieda}


\item 已知关于$x$的方程$\sin^2x+\cos{x}+m=0$. 
\begin{enumerate}
\item 当$m=1$时, 求方程的解; 
\item 要使此方程有解, 试确定$m$的取值范围. 
\end{enumerate}
\begin{jieda}
\begin{enumerate}[(1)]
\item $1-\cos^2x+\cos{x}+1=0$, 解得$\cos{x}=-1$(2已舍去), 故$x=2k\pi+\pi(k\zz)$. 

\item $m=-\sin^2x-\cos{x}=\cos^2x-\cos{x}-1\in[-\frac{5}{4},1]$. 
\end{enumerate}
\end{jieda}
\end{enumerate}





\subsection{正弦、余弦、正切、余切阶段练习}

\begin{enumerate}
\item 已知$\sin\alpha+\cos\alpha=\dfrac{1}{5}$, $\alpha$是第二象限角, 那么$\tan\alpha$=\blkda{$-\dfrac{4}{3}$}. 

\item 已知$\cot\theta=-\dfrac{3}{4}$, 则$\cos\theta$=\blkda{$\pm\dfrac{3}{5}$}. 

\item 若$\tan{x}+\cot{x}=-2$, 则$\sin{x}+\cos{x}$=\blkda{0}. 

\item 已知$\cos\left(\dfrac{\pi}{6}-\alpha\right)=\dfrac{\sqrt{3}}{3}$, 
则$\sin^2\left(\alpha-\dfrac{\pi}{6}\right)-\cos\left(\dfrac{5\pi}{6}+\alpha\right)$的值为\blkda{$\dfrac{2+\sqrt{3}}{3}$}. 

\item 已知$\sin(\pi-\alpha)+2\cos(\pi+\alpha)=0$, 则$\dfrac{1}{\sin\alpha\cos\alpha}$=\blkda{$\dfrac{5}{2}$}. 
\begin{jieda}
$\tan\alpha=2$, $E=\dfrac{1+\tan^2\alpha}{\tan\alpha}=\dfrac{5}{2}$
\end{jieda}

\item 化简: $(\sec^2\alpha+\csc^2\alpha)-(\cot\alpha-\tan\alpha)^2$=\blkda{4}. 

\item 若$\sin\alpha+\sin^2\alpha=1$, 则$\cos^{2}\alpha+\cos^{4}\alpha$的值等于\blkda{1}. 
\begin{jieda}
$\iff\sin\alpha=\cos^2\alpha$, 故$E=\cos^2\alpha+\sin^2\alpha=1$. 
\end{jieda}

\item 若$0<\alpha<\dfrac{\pi}{2}$, 且$\lg(1+\cos\alpha)=m$, $\lg\dfrac{1}{1-\cos\alpha}=n$, 则$\lg\sin\alpha$=\blkda{$\dfrac{m-n}{2}$}(用$m,n$表示). 
\begin{jieda}
因为$\alpha$为锐角, 故$\sin\alpha=\sqrt{1-\cos^2\alpha}$. \\
于是$\lg\sin\alpha=\frac{1}{2}\lg(1-\cos^2\alpha)=\frac{1}{2}(\lg(1+\cos\alpha)+\lg(1-\cos\alpha))=\frac{m-n}{2}$. 
\end{jieda}

\item ``$\sin{x}=\dfrac{1}{2}$''是``$x=2k\pi+\dfrac{\pi}{6}(k\zz)$''的\dotfill(\kh{B})
\xx{充分不必要条件}{必要不充分条件}{充要条件}{既不充分也不必要条件}

\item 若$\sqrt{\dfrac{1-\sin{x}}{1+\sin{x}}}=\dfrac{\sin{x}-1}{\cos{x}}$, 则$x$的取值范围是\dotfill(\kh{B})
\xx{$2k\pi\leq x\leq 2k\pi+\dfrac{\pi}{2}(k\zz)$}
{$2k\pi+\dfrac{\pi}{2}<x<2k\pi+\dfrac{3\pi}{2}(k\zz)$}
{$2k\pi+\dfrac{3\pi}{2}<x<2(k+1)\pi(k\zz)$}
{$(2k+1)\pi<x<2k\pi+\dfrac{3\pi}{2}(k\zz)$}

\item 若$2k\pi-\dfrac{\pi}{4}\leq \alpha\leq 2k\pi+\dfrac{\pi}{4}(k\zz)$, 
化简$\sqrt{1-2\sin\alpha\cos\alpha}+\sqrt{1+2\sin\alpha\cos\alpha}$的结果是\dotfill(\kh{C})
\xx{$2\sin\alpha$}{$-2\sin\alpha$}{$2\cos\alpha$}{$-2\cos\alpha$}
\begin{jieda}
所求非负, 于是看选项只有C符合. 
\end{jieda}

\item 若$\sin\alpha-\csc\alpha<0$, 则角$\alpha$所在象限是\dotfill(\kh{A})
\xx{第一、二象限}{第二、三象限}{第三、四象限}{第二、四象限}
\begin{jieda}
$\iff\sin\alpha-\frac{1}{\sin\alpha}<0\iff\sin\alpha\in(0,1)$. 
\end{jieda}


\item 已知$f(x)=2\cos\dfrac{\pi x}{6}$, 则$f(0)+f(1)+f(2)+\cdots+f(2021)=$\dotfill(\kh{C})
\xx{$-\sqrt{3}$}{$\sqrt{3}$}{2}{$-2$}
% \xx{$-\sqrt{3}$}{$\sqrt{3}$}{1}{$-1$}  %  原稿中的. 20230103
\begin{jieda}
$\sum\limits_{k=l+1}^{l+12}f(k)=0$, $2021=12\cdot168+5$, 
故$E=f(0)+f(1)+f(2)+f(3)+f(4)+f(5)=f(0)=2$. \\
\ps 以后修改为填空题吧. 每次改数字就是不修改选项. (选项已修改)
\end{jieda}

\item 若角$A$,$B$,$C$是$\triangle{ABC}$的三个内角, 则下列等式中一定成立的是\dotfill(\kh{D})
\xx{$\cos(A+B)=\cos{C}$}
{$\sin(A+B)=-\sin{C}$}
{$\cos\left(\dfrac{A}{2}+C\right)=\sin{B}$}
{$\sin\dfrac{B+C}{2}=\cos\dfrac{A}{2}$}

\item $\triangle{ABC}$的内角$A$满足$\sin{A}+\cos{A}>0$且$\tan{A}-\sin{A}<0$, 则$A$的取值范围是\dotfill(\kh{B})
\xx{$\left(0,\dfrac{\pi}{4}\right)$}
{$\left(\dfrac{\pi}{2},\dfrac{3\pi}{4}\right)$}
{$\left(\dfrac{\pi}{4},\dfrac{\pi}{2}\right)$}
{$\left(\dfrac{\pi}{4},\dfrac{3\pi}{4}\right)$}
\begin{jieda}
单位圆中: $\sin{A}+\cos{A}>0\iff y+x>0$, 于是$A\in(0,\frac{3\pi}{4})$. \\
又$\tan{A}<\sin{A}\iff\dfrac{1}{\cos{A}}<1$, 故$A$钝角. 
\end{jieda}

\item 设$f(x)=(\sin\alpha)^{x-\frac{1}{2}}$, 若$f(\log_2\sin\alpha)=\sqrt{2}$, 且$0<\alpha<\dfrac{\pi}{2}$, 则$\alpha$的值是\dotfill(\kh{A})
\xx{$\dfrac{\pi}{4}$}{$\dfrac{\pi}{6}$}{$\dfrac{\pi}{3}$}{$\dfrac{\pi}{8}$}
\begin{jieda}
令$t=\sin\alpha$, 则$t^{\log_2t-\frac{1}{2}}=\sqrt{2}$, 故$(\log_2t-\frac{1}{2})\log_2t=\log_2\sqrt{2}=\frac{1}{2}$. \\
解得$\log_2t=1$或$-\frac{1}{2}$, 故$\sin\alpha=2$或$\frac{\sqrt{2}}{2}$. \\
又$\alpha$为锐角, 故$\alpha=\frac{\pi}{4}$. 
\end{jieda}

\item 已知角$\alpha$的终边过点$A(-1,m)$, 且$\sin\alpha=\dfrac{\sqrt{5}}{5}m(m\neq0)$. 
\begin{enumerate}[(1)]
\item 求非零实数$m$的值;
\item 当$m>0$时, 求$\dfrac{\sin(2\pi-\alpha)+\cos(\pi+\alpha)}{\cos(\alpha-\pi)-\cos(\frac{3\pi}{2}-\alpha)}$的值. 
\end{enumerate}
\begin{jieda}
\begin{enumerate}[(1)]
\item $\dfrac{\sqrt{5}}{5}m=\dfrac{m}{\sqrt{1+m^2}}\land m\neq 0$, 解得$m=\pm2$. 
\item 因为$m>0$, 故$m=2$. 终边过点$A(-1,2)$, 故$\tan\alpha=-2$. \\
$E=\dfrac{-\sin\alpha-\cos\alpha}{-\cos\alpha+\sin\alpha}=\dfrac{\tan\alpha+1}{-\tan\alpha+1}=-\dfrac{1}{3}$. 
\end{enumerate}
\end{jieda}


\item
\begin{enumerate}[(1)]
\item 已知$\tan\alpha-\cot\alpha=\dfrac{1}{4}$, 求$\tan^3\alpha-\cot^3\alpha$的值; 
\item 已知$\dfrac{\sin\alpha+\cos\alpha}{\sin\alpha-\cos\alpha}=2$, 求$\sin\alpha\cos\alpha$的值. 
\end{enumerate}
\begin{jieda}
\begin{enumerate}[(1)]
\item 令$t=\tan\alpha$, 则$t-\frac{1}{t}=\frac{1}{4}$. \\
$E=t^3-\frac{1}{t^3}=(t-\frac{1}{t})(t^2+1+\frac{1}{t^2})=(t-\frac{1}{t})((t-\frac{1}{t})^2+3)=\frac{49}{64}$. 

\item $\iff\sin\alpha+\cos\alpha=2\sin\alpha-2\cos\alpha\iff\tan\alpha=3$. \\
$E=\dfrac{\sin\alpha\cos\alpha}{\sin^2\alpha+\cos^2\alpha}=\dfrac{\tan\alpha}{\tan^2\alpha+1}=\dfrac{3}{10}$. 
\end{enumerate}
\end{jieda}

\item
\begin{enumerate}[(1)]
\item 已知$\cos\theta-\sin\theta=\sqrt{2}\sin\theta$, 求证: $\tan\theta=\dfrac{\cos\theta-\sin\theta}{\cos\theta+\sin\theta}$;
\item 已知$1-a\sin\alpha=\cos\alpha$, $1+\cos\alpha=b\sin\alpha$, 求证: $ab=1$; 
\item 已知$\begin{cases}
\tan\alpha+\sin\alpha=m\\
\tan\alpha-\sin\alpha=n
\end{cases}$, 求证: $(m^2-n^2)^2=16mn$. 
\end{enumerate}
\begin{jieda}
\begin{enumerate}[(1)]
\item $\tan\theta=\sqrt{2}-1$. 
\item $\begin{cases}
1-\cos\alpha=a\sin\alpha\\
1+\cos\alpha=b\sin\alpha
\end{cases}$, 两式相乘即得. 
\item 解出$\tan\alpha,\sin\alpha$, 然后$1+\tan^2\alpha=\sec^2\alpha=\dfrac{1}{\cos^2\alpha}=\dfrac{1}{1-\sin^2\alpha}$. 总归能代换掉. 
\end{enumerate}
\end{jieda}

\item 求函数$y=\dfrac{a^2}{\cos^2x}+\dfrac{b^2}{\sin^2x}(a,b>0,0<x<\frac{\pi}{2})$的最小值, 并求取到最小值时$x$的正切值(用$a,b$表示). 
\begin{jieda}
分式型柯西不等式(权方和的特殊情形), 或者用正切表示后用基本不等式. \\
柯西: $y\geq\dfrac{(a+b)^2}{\cos^2x+\sin^2x}=(a+b)^2$. \\
正切: $y=\dfrac{a^2(\sin^2x+\cos^2x)}{\cos^2x}+\dfrac{b^2(\sin^2x+\cos^2x)}{\sin^2x}=a^2\tan^2x+a^2+b^2+\dfrac{b^2}{\tan^2x}\geq a^2+b^2+2ab$. \\
\da $\min=(a+b)^2$, $\tan{x}=\sqrt{\dfrac{b}{a}}$. 
\end{jieda}

\item 一个大风车的半径为8\si{m}, 12分钟旋转一周, 它的最低点离地面2\si{m}(如图所示), 求风车翼片的一个端点离地面距离$h$(米)
与时间$t$(分钟)之间的函数关系(用弧度制求解). \\
\begin{minipage}{\textwidth}
\centering
\includegraphics[totalheight=4cm]{./bodyxb/src/main42chap22sanjiaobizonghe01.png}
\end{minipage}
\begin{jieda}
坐标原点选在圆心, $x$轴为水平方向向右, $y$轴为竖直方向向上. 设$P$对应的角的大小为$\theta$. \\
若$t=0$时, 点$P$坐标为$(8,0)$, 则$\theta=\omega t=\dfrac{2\pi t}{12}$. \\
于是$h=10+8\sin\theta=10+8\sin\frac{\pi t}{6}$. \\
\ps 坐标系如何建立得说清楚, 特别是$\theta$具体所指要在图中标出, 且在解答过程中加以说明. 
\end{jieda}



\item (选做题)已知$\alpha$为锐角, 
\begin{enumerate}[(1)]
\item 试证: $1<\sin\alpha+\cos\alpha\leq\sqrt{2}$; 
\item 探究$\alpha,\sin\alpha,\tan\alpha$之间的大小关系. 
\end{enumerate}
\begin{jieda}
说是三角函数线, 也可以不是. 
\begin{enumerate}[(1)]
\item 单位圆中看. 
\item 弧度制条件下, $\sin\alpha<\alpha<\tan\alpha$. 
\end{enumerate}
\end{jieda}
\end{enumerate}






\section{常用三角公式}

\subsection{两角和与差的余弦、正弦和正切公式(1)}
\zsd 两角和与差的余弦公式

\begin{enumerate}
\item 若$\alpha$是第一象限角, $\sin\alpha=\dfrac{1}{4}$, 则$\cos\left(\alpha+\frac{\pi}{3}\right)$=\blkda[2]{$\dfrac{\sqrt{15}-\sqrt{3}}{8}$}. 

\item 若$\triangle{ABC}$中, $\cos{A}=-\dfrac{5}{13}$, $\sin{B}=\dfrac{4}{5}$, 则$\cos{C}$=\blkda{$\dfrac{63}{65}$}. 
\begin{jieda}
$\cos{B}=\pm\frac{3}{5}$, 则$-\frac{3}{5}$如何舍去? \\
当前情形下, 用$\sin{C}$辅助. 
\end{jieda}

\item 化简: $\cos\left(\dfrac{\pi}{2}+\alpha\right)\tan\left(\dfrac{3\pi}{2}-\alpha\right)$=\blkda{$-\cos\alpha$}. 

\item 若$\triangle{ABC}$中, 若$\cos{A}\cos{B}>\sin{A}\sin{B}$, 则$\triangle{ABC}$是\blkda{钝角}三角形. 

\item 已知$\cos(\alpha-\beta)=\dfrac{5}{13}$, $\sin\alpha=-\dfrac{4}{5}$, $\alpha\in\left(-\dfrac{\pi}{2},\dfrac{\pi}{2}\right)$, 
$\beta\in(0,\pi)$, 则$\cos\beta$=\blkda{$\dfrac{63}{65}$}. 
\begin{jieda}
$\sin\alpha=-\dfrac{4}{5}$, $\alpha\in\left(-\dfrac{\pi}{2},\dfrac{\pi}{2}\right)$, 则$\alpha\in\left(-\dfrac{\pi}{2},0\right)$, 且$\cos\alpha=\dfrac{3}{5}$. \\
于是$\alpha-\beta\in\left(-\dfrac{3\pi}{2},0\right)$. \\
又$\cos(\alpha-\beta)=\dfrac{5}{13}>0$, 故$\alpha-\beta\in\left(-\dfrac{\pi}{2},0\right)$, 且$\sin(\alpha-\beta)=-\dfrac{12}{13}$. \\
于是$\cos\beta=\cos(-\beta)=\cos(\alpha-\beta-\alpha)=\dfrac{63}{65}$. \\
\ps 三角比求值, 麻烦在于确定角的范围. 象限角与三角比的符号相结合是常见的判断角度. 
\end{jieda}


\item 化简: $\cos\alpha+\cos(120^{\circ}-\alpha)+\cos(120^{\circ}+\alpha)$=\blkda{0}. 

\item 若$\sin2x\sin3x=\cos2x\cos3x$, 则$x$的一个值是\dotfill(\kh{A})
\xx{$18^{\circ}$}{$36^{\circ}$}{$45^{\circ}$}{$30^{\circ}$}

\item 化简: $\dfrac{\sin(\pi-x)\cos(\frac{3\pi}{2}-x)\tan{x}}{\cos(2\pi-x)\cot(x-\frac{\pi}{2})\sin(\frac{\pi}{2}-x)}=$\dotfill(\kh{B})
\xx{$-1$}{$\tan^2x$}{1}{$-\tan^2x$}

\item 已知$\sin\alpha=\dfrac{2}{3}$, $\cos\beta=-\dfrac{3}{4}$, 且两角在同一个象限内, 则$\cos(\alpha-\beta)$的值为\dotfill(\kh{A})
\xx{$\dfrac{3\sqrt5+2\sqrt7}{12}$}
{$\dfrac{3\sqrt5-2\sqrt7}{12}$}
{$\dfrac{-3\sqrt5+2\sqrt7}{12}$}
{$-\dfrac{3\sqrt5+2\sqrt7}{12}$}
\begin{jieda}
因为两角在同一个象限内, 且$\sin\alpha>0$, $\cos\beta<0$, 故$\alpha,\beta$都是第二象限角. \\
于是$\cos\alpha=-\dfrac{\sqrt{5}}{3}$, $\sin\beta=\dfrac{\sqrt{7}}{4}$, 故$\cos(\alpha-\beta)=\dfrac{3\sqrt{5}+2\sqrt{7}}{12}$. 
\end{jieda}


\item 已知$\cos\alpha=-\dfrac{8}{17}$, $\cos\beta=-\dfrac{12}{13}$, 且两角为不同象限的角, 则$\cos(\alpha-\beta)$的值为\dotfill(\kh{D})
\xx{$\dfrac{171}{221}$}
{$-\dfrac{21}{221}$}
{$\pm\dfrac{21}{221}$}
{$\dfrac{21}{221}$}
\begin{jieda}
因为两角为不同象限的角, 且余弦都小于零, \\
故两者一为第二象限角, 一为第三象限角, 于是$\sin\alpha\sin\beta<0$. \\
故$\cos(\alpha-\beta)=\cos\alpha\cos\beta+\sin\alpha\sin\beta=\dfrac{8\cdot12-15\cdot5}{17\cdot13}=\dfrac{21}{221}$. 
\end{jieda}


\item 已知$\cos(\alpha+\beta)=\dfrac{4}{5}$, $\cos(\alpha-\beta)=-\dfrac{5}{13}$, 
$\dfrac{3\pi}{2}<\alpha+\beta<2\pi$, $\dfrac{\pi}{2}<\alpha-\beta<\pi$, 求$\cos2\alpha$和$\cos2\beta$的值. 
\begin{jieda}
结合角的范围, 易得$\sin(\alpha+\beta)=-\dfrac{3}{5}$, $\sin(\alpha-\beta)=\dfrac{12}{13}$. \\
$\cos2\alpha=\cos(\alpha+\beta+\alpha-\beta)=\dfrac{16}{65}$. \\
$\cos2\beta=\cos(\alpha+\beta-(\alpha-\beta))=-\dfrac{56}{65}$. 
\end{jieda}


\item 已知$\cos(2\alpha-\beta)=-\dfrac{11}{14}$, $\sin(\alpha-2\beta)=\dfrac{4\sqrt{3}}{7}$, 
$0<\beta<\dfrac{\pi}{4}<\alpha<\dfrac{\pi}{2}$, 求$\alpha+\beta$的值. 
\begin{jieda}
$\dfrac{\pi}{4}<2\alpha-\beta<\pi$, 故$\sin(2\alpha-\beta)=\dfrac{5\sqrt{3}}{14}$. \\
$-\dfrac{\pi}{4}<\alpha-2\beta<\dfrac{\pi}{2}$, 故$\cos(\alpha-2\beta)=\dfrac{1}{7}$. \\
于是$\cos(\alpha+\beta)=\cos((2\alpha-\beta)-(\alpha-2\beta))=\dfrac{1}{2}$, 且$\alpha+\beta\in\left(\dfrac{\pi}{4},\dfrac{3\pi}{4}\right)$, \\
故$\alpha+\beta=\dfrac{\pi}{3}$. 
\end{jieda}


\item 若$\sin\alpha+\sin\beta+\sin\gamma=0$, $\cos\alpha+\cos\beta+\cos\gamma=0$, 求$\cos(\alpha-\beta)$. 
\begin{jieda}
$1=\sin^2\gamma+\cos^2\gamma=(\sin\alpha+\sin\beta)^2+(\cos\alpha+\cos\beta)^2=2+2\cos(\alpha-\beta)$, \\
展开化简易得$\cos(\alpha-\beta)=-\dfrac{1}{2}$. \\
\ps 几何上, $\alpha,\beta,\gamma$终边之间的夹角为$\ang{120}$. 
\end{jieda}

\item 已知$8\cos(2\alpha+\beta)+5\cos\beta=0$, 求$\tan(\alpha+\beta)\cdot\tan\alpha$的值. 
\begin{jieda}
$\iff 8\cos((\alpha+\beta)+\alpha)+5\cos((\alpha+\beta)-\alpha)=0\iff 13\cos(\alpha+\beta)\cos\alpha=3\sin(\alpha+\beta)\sin\alpha$. \\
(1) 若$\cos\alpha=0$, 则$\sin\alpha=\pm1\neq 0$, 于是$\sin(\alpha+\beta)=0$, \\
故$0=\sin\alpha\cos\beta+\cos\alpha\sin\beta$, 故$\cos\beta=0$. \\
此时$\tan(\alpha+\beta)\cdot\tan\alpha$无意义. \\
(2) 若$\cos(\alpha+\beta)=0$, 则$\sin(\alpha+\beta)=\pm1\neq0$, 于是$\sin\alpha=0$. \\
而$0=\cos(\alpha+\beta)=\cos\alpha\cos\beta-\sin\alpha\sin\beta$, 故$\cos\beta=0$. \\
此时$\tan(\alpha+\beta)\cdot\tan\alpha$无意义. \\
% 于是$\alpha=m\pi+\dfrac{\pi}{2},\alpha+\beta=n\pi$, 
% 即$\alpha=m\pi+\dfrac{\pi}{2},\beta=n\pi+\dfrac{\pi}{2}$. 
(3) 若$\cos(\alpha+\beta)\cos\alpha\neq 0$, 则$\tan(\alpha+\beta)\tan\alpha=\dfrac{13}{3}$. \\
综上, 当$\alpha=\frac{m\pi}{2}$, $n=n\pi+\frac{\pi}{2}$时, $\tan(\alpha+\beta)\cdot\tan\alpha$无意义; \\
当$\alpha\neq\frac{m\pi}{2}$, $n\neq n\pi+\frac{\pi}{2}$时, $\tan(\alpha+\beta)\tan\alpha=\dfrac{13}{3}$. 
\end{jieda}

\end{enumerate}







\subsection{两角和与差的余弦、正弦和正切公式(2)}
\zsd 两角和与差的正弦公式

\begin{enumerate}
\item 化简$\sin(x+27^\circ)\cos(18^\circ-x)+\sin(18^\circ-x)\cos(x+27^\circ)$=\blkda{$\dfrac{\sqrt{2}}{2}$}. 

\item 已知$\cos\alpha=-\dfrac{1}{2}$, $\sin\beta=-\dfrac{\sqrt{3}}{2}$, 且$\alpha\in\left(\dfrac{\pi}{2},\pi\right)$, $\beta\in\left(\dfrac{3\pi}{2},2\pi\right)$,
则$\sin(\alpha+\beta)$=\blkda{$\dfrac{\sqrt{3}}{2}$}. 

\item 若$\alpha,\beta$均为锐角, $P=\sin(\alpha+\beta)$, $Q=\sin\alpha+\sin\beta$, 则$P$与$Q$的大小关系是\blkda{$P<Q$}. 
\begin{jieda}
法一: $\alpha,\beta$均为锐角, 故$\cos\alpha,\sin\alpha,\cos\beta,\sin\beta\in(0,1)$. \\
于是$P=\sin\alpha\cos\beta+\cos\alpha\sin\beta<\sin\alpha\cdot 1+1\cdot\sin\beta=Q$, 即$P<Q$. \\
法二: $P=2\sin\frac{\alpha+\beta}{2}\cos\frac{\alpha+\beta}{2}$, $Q=2\sin\frac{\alpha+\beta}{2}\cos\frac{\alpha-\beta}{2}$. 
\end{jieda}

\item 求值: $\dfrac{\sin22^{\circ}+\cos30^{\circ}\sin8^{\circ}}{\cos22^{\circ}-\sin30^{\circ}\sin8^{\circ}}$=\blkda{$\dfrac{\sqrt{3}}{3}$}. 

\item 命题甲: $3\sin\beta=\sin(2\alpha+\beta)$是命题乙: $2\tan\alpha=\tan(\alpha+\beta)$成立的\blkda[2.5]{必要不充分}条件. 
\begin{jieda}
甲$\iff3\sin(\alpha+\beta-\alpha)=\sin(\alpha+\beta+\alpha)\iff \sin(\alpha+\beta)\cos\alpha=2\cos(\alpha+\beta)\sin\alpha$. \\
若$\cos\alpha=0$, 则$\sin\alpha\neq 0$, 故$\cos(\alpha+\beta)=0$; 同理, 若$\cos(\alpha+\beta)=0$, 则$\cos\alpha=0$. \\
此时$\tan\alpha,\tan(\alpha+\beta)$无意义($\alpha=m\pi+\frac{\pi}{2},\beta=n\pi(m,n\zz)$). \\
若$\cos\alpha\cos(\alpha+\beta)\neq 0$, 则易得乙成立. \\
综上, 甲是乙的``必要不充分''条件. 
\end{jieda}


\item 下面这道填空题因印刷原因造成在横线上的内容无法认清, 现知结论, 请在横线上填写原题的一个条件. \\
题目: 已知$\alpha,\beta$均为锐角, 且$\sin\alpha-\cos\beta=-\dfrac{1}{2}$,\blkda[3]{$\cos\alpha+\sin\beta=\dfrac{1}{3}$}, 
则$\sin(\alpha-\beta)=\dfrac{59}{72}$. 
\begin{jieda}
设$\cos\alpha+\sin\beta=p$, 则$p^2+\dfrac{1}{4}=2-2\sin(\alpha-\beta)$, 故$p^2=\dfrac{1}{9}$. \\
又$\alpha,\beta$均为锐角, 故$p=\dfrac{1}{3}$. \\
\ps 从上述过程可以, $\cos\alpha+\sin\beta$必须是$\dfrac{1}{3}$. 由于无法验证所有情形, 故所填内容条件必须与上述结果相容. 
\end{jieda}


\item 已知$\sin\alpha\cos60^{\circ}-\cos\alpha\sin60^{\circ}=\frac{1}{2}$, $\alpha\in(0,2\pi)$, 则$\alpha=$\dotfill(\kh{D})
\xx
{$\dfrac{\pi}{2}$}
{$\dfrac{7\pi}{6}$}
{$\dfrac{\pi}{6}$或$\dfrac{3\pi}{2}$}
{$\dfrac{\pi}{2}$或$\dfrac{7\pi}{6}$}

\item 若$\sin(\alpha+\beta)\sin(\beta-\alpha)=m$, 则$\sin^2\beta-\sin^2\alpha$为\dotfill(\kh{B})
\xx{$-m$}{$m$}{$-4m$}{$4m$}
\begin{jieda}
\begin{align*}
\sin(\alpha+\beta)\sin(\alpha-\beta) 
& =\sin^2\alpha\cos^2\beta-\cos^2\alpha\sin^2\beta\\
& =\sin^2\alpha(1-\sin^2\beta)-(1-\sin^2\alpha)\sin^2\beta\\
& =\sin^2\alpha-\sin^2\beta\\
& =\cos^2\beta-\cos^2\alpha
\end{align*}
同理$\cos(\alpha+\beta)\cos(\alpha-\beta)=\cos^2\alpha-\sin^2\beta=\cos^2\beta-\sin^2\alpha$. 
\end{jieda}


\item 已知$\triangle{ABC}$中, 已知$\sin{A}\sin{B}+\sin{A}\cos{B}+\cos{A}\sin{B}+\cos{A}\cos{B}=2$, 则$\triangle{ABC}$是\dotfill(\kh{C})
\xx{等腰三角形}{直角三角形}{等腰直角三角形}{等腰或直角三角形}
\begin{jieda}
$\iff 2=\cos(A-B)+\sin(A+B)\iff 1=\cos(A-B)=\sin(A+B)$. 
\end{jieda}


\item 下列四个命题中, 假命题是\dotfill(\kh{B})
\xx
{存在这样的$\alpha,\beta$值, 使得$\sin(\alpha+\beta)=\sin\alpha\cos\beta-\cos\alpha\sin\beta$}
{不存在无穷多的$\alpha,\beta$值, 使得$\sin(\alpha+\beta)=\sin\alpha\cos\beta-\cos\alpha\sin\beta$}
{对于任意的$\alpha,\beta$值, 都有$\sin(\alpha+\beta)=\sin\alpha\cos\beta+\cos\alpha\sin\beta$}
{不存在$\alpha,\beta$值, 使得$\sin(\alpha+\beta)\neq\sin\alpha\cos\beta+\cos\alpha\sin\beta$}
\begin{jieda}
对于B选项, 只要$\cos\alpha\sin\beta=0$, 则必然有$\sin(\alpha+\beta)=\sin\alpha\cos\beta-\cos\alpha\sin\beta$. \\
C与D等价. 
\end{jieda}

\item 化简: 
\begin{enumerate}[(1)]
\item $\sin260^\circ\cos40^\circ-\cos(-260^\circ)\cos130^\circ$;
\item $\cos(15^\circ-\alpha)\sin15^\circ-\sin(165^\circ+\alpha)\cos(-15^\circ)$. 
\end{enumerate}
\begin{jieda}
\begin{enumerate}[(1)]
\item $E=-\sin\ang{80}\cos\ang{40}-\cos\ang{80}\sin\ang{40}=-\sin\ang{120}=-\dfrac{\sqrt{3}}{2}$. 
\item $E=\cos(\ang{15}-\alpha)\sin\ang{15}-\sin(\ang{15}-\alpha)\cos\ang{15}=\sin(\ang{15}-(\ang{15}-\alpha))=\sin\alpha$. 
\end{enumerate}
\end{jieda}

\item 已知$2\sin\beta=\sin(2\alpha+\beta)$, 求$\tan(\alpha+\beta):\tan\alpha$的值. 
\begin{jieda}
$\iff 2\sin(\alpha+\beta-\alpha)=\sin(\alpha+\beta+\alpha)\iff\sin(\alpha+\beta)\cos\alpha=3\cos(\alpha+\beta)\sin\alpha$. \\
若$\cos\alpha=0$, 则$\cos(\alpha+\beta)=0$; 若$\cos(\alpha+\beta)=0$, 则$\cos\alpha=0$. \\
故当$\cos\alpha\cos(\alpha+\beta)=0$时, $\tan(\alpha+\beta):\tan\alpha$不存在; \\
当$\cos\alpha\cos(\alpha+\beta)\neq 0$时, $\tan(\alpha+\beta):\tan\alpha=3:1$. 
\end{jieda}


\item 已知$0^\circ<\alpha<\beta<90^\circ$, 且$\sin\alpha,\sin\beta$是方程$x^2-(\sqrt2\cos40^\circ)x+\cos^240^{\circ}-\frac{1}{2}=0$的两个实根, 
求$\sin(\beta-5\alpha)$和$\sin(\beta-2\alpha)$的值. 
\begin{jieda}
法一: 方程的解为$x=\dfrac{\sqrt{2}}{2}(\cos\ang{40}\pm\sin\ang{40})$, 即$x=\sin\ang{85}$或$\sin\ang{5}$. \\
于是结合条件可得$\alpha=\ang{5},\beta=\ang{85}$. \\
法二: $\begin{cases}
\sin\alpha+\sin\beta=\sqrt{2}\cos\ang{40}\\
\sin\alpha\sin\beta=\cos^2\ang{40}-\frac{1}{2}
\end{cases}$, \\
于是$(\sin\alpha+\sin\beta)^2=2\sin\alpha\sin\beta+1$, 故$\sin^2\alpha+\sin^2\beta=1$. \\
于是$\sin^2\alpha=\cos^2\beta$. 又$0^\circ<\alpha<\beta<90^\circ$, 故$\sin\alpha=\cos\beta$, 即$\alpha+\beta=\ang{90}$. \\
代入第一式可得, $\sin\alpha+\cos\alpha=\sqrt{2}\sin(\alpha+\ang{45})=\sqrt{2}\cos\ang{40}$, 即$\alpha+\ang{45}=\ang{50},\ang{130}$. \\
结合条件可得$\alpha=\ang{5}$, $\beta=\ang{85}$. \\
于是$\sin(\beta-5\alpha)=\dfrac{\sqrt{3}}{2}$, $\sin(\beta-2\alpha)=\dfrac{\sqrt{6}+\sqrt{2}}{4}$. \\
\ps 由方程可得$\cos\ang{40}=\dfrac{\sqrt{2}}{2}(x\pm\sqrt{1-x^2})$, 一般化为$\sqrt{2}\cos\theta=x\pm\sqrt{1-x^2}$, 对照上述过程可知$x=\sin(\ang{45}\pm\theta)$. \\
法三: $(\sin\beta-\sin\alpha)^2=(\sin\beta+\sin\alpha)^2-4\sin\alpha\sin\beta=2\sin^2\ang{40}$, 于是$\sin\beta-\sin\alpha=\sqrt{2}\sin\ang{40}$. \\
于是$\sin\alpha=\sin\ang{5}$, $\sin\beta=\sin\ang{85}$. 
\end{jieda}

\item 已知$\sin2(\alpha+\gamma)=n\sin2\beta$, 求$\dfrac{\tan(\alpha+\beta+\gamma)}{\tan(\alpha-\beta+\gamma)}$的值. 
\begin{jieda}
法一: $\iff\sin((\alpha+\beta+\gamma)+(\alpha-\beta+\gamma))=n\sin((\alpha+\beta+\gamma)-(\alpha-\beta+\gamma))$. \\
展开合并化简易得$\dfrac{\tan(\alpha+\beta+\gamma)}{\tan(\alpha-\beta+\gamma)}=\dfrac{n+1}{n-1}$. \\
法二: $\dfrac{\tan{x}+\tan{y}}{\tan{x}-\tan{y}}=\dfrac{\sin{x}\cos{y}+\cos{x}\sin{y}}{\sin{x}\cos{y}-\cos{x}\sin{y}}=\dfrac{\sin(x+y)}{\sin(x-y)}$. \\
令$a=\tan(\alpha+\beta+\gamma), b=\tan(\alpha-\beta+\gamma)$, \\
则$\dfrac{a+b}{a-b}=\dfrac{\sin2(\alpha+\gamma)}{\sin2\beta}=n$, 
故$\dfrac{a}{b}=\dfrac{n+1}{n-1}$. 
\end{jieda}
\end{enumerate}





\subsection{两角和与差的余弦、正弦和正切公式(3)}
\zsd 两角和与差的正切公式

\begin{enumerate}
\item $\tan\alpha=-2$, 则$\tan\left(\alpha-\frac{\pi}{3}\right)$=\blkda{$\dfrac{8+5\sqrt{3}}{11}$}. 

\item $\tan\alpha=-3$, $\tan(\alpha-\beta)=\frac{1}{2}$, 则$\tan\beta$=\blkda{7}.  

\item 化简: $\dfrac{\tan(20^{\circ}+\alpha)+\tan(40^{\circ}-\alpha)}{1-\tan(20^{\circ}+\alpha)\tan(40^{\circ}-\alpha)}$=\blkda{$\sqrt{3}$}. 

\item 已知$\tan\alpha=\dfrac{1}{7}$, $\tan\beta=\dfrac{1}{3}$, $\alpha,\beta$均为锐角, 则$\alpha+2\beta$=\blkda{$\dfrac{\pi}{4}$}.

\item 已知$\tan(\alpha+\beta)=\dfrac{2}{5}$, $\tan\left(\beta-\dfrac{\pi}{4}\right)=\dfrac{1}{4}$, 则$\tan\left(\alpha+\dfrac{\pi}{4}\right)$=\blkda{$\dfrac{3}{22}$}. 

\item 已知$\alpha$为第三象限角, $\sin\alpha=-\dfrac{\sqrt3}{3}$, 则$\dfrac{\tan(\alpha+\beta)-\tan\beta}{1+\tan(\alpha+\beta)\tan\beta}$=\blkda{$\dfrac{\sqrt{2}}{2}$}. 

\item 若$\tan\alpha=2$, 则$\dfrac{1-\tan(\alpha-45^{\circ})}{1+\tan(\alpha-45^{\circ})}$的值为\dotfill(\kh{B})
\xx{1}{$\dfrac{1}{2}$}{$\dfrac{\sqrt3}{3}$}{$\sqrt3$}
\begin{jieda}
$\dfrac{1-\tan(\alpha-45^{\circ})}{1+\tan(\alpha-45^{\circ})}=\tan(\ang{45}-(\alpha-\ang{45}))=\cot\alpha=\dfrac{1}{2}$. 
\end{jieda}


\item 若$\tan\alpha,\tan\beta$是方程$x^2-px+q=0$的两根, $\cot\alpha,\cot\beta$是方程$x^2-rx+s=0$的两根, 
则 (1) $ps=r$; (2) $qs=1$; (3) $qr=p$; (4) $r(1-q)=p(s-1)$中成立的式子有几个\dotfill(\kh{D})
\xx{1}{2}{3}{4}
\begin{jieda}
$\tan\alpha\cot\alpha=\tan\beta\cot\beta=1$. \\
于是方程$x^2-px+q=0$与$sx^2-rx+1=0$是同解方程, 且$s\neq 0$, $\Delta\geq 0$. \\
于是$\dfrac{1}{s}=\dfrac{-p}{-r}=\dfrac{q}{1}$, 即$\dfrac{1}{s}=\dfrac{p}{r}=q$. (若分母为零, 则改写为乘积的形式)\\
易得(1)(2)(3)成立, 将(1)(2)(3)代入(4)中(4)也成立. 
\end{jieda}


\item 若$\sin\alpha\sin\beta+\cos\alpha\cos\beta=0$, 那么$\sin\alpha\cos\alpha+\sin\beta\cos\beta$的值等于\dotfill(\kh{B})
\xx{$-1$}{0}{$\dfrac{\sqrt2}{2}$}{1}
\begin{jieda}
$\iff\cos(\alpha-\beta)=0\iff\alpha-\beta=k\pi+\dfrac{\pi}{2}$. \\
于是$E=\sin(\beta+k\pi+\frac{\pi}{2})\cos(\beta+k\pi+\frac{\pi}{2})+\sin\beta\cos\beta=0$. (对$k$分奇偶讨论一下)
\end{jieda}

\item 化简$\dfrac{\tan(x+y)-\tan{x}-\tan{y}}{\tan{x}\tan(x+y)}$得到的结果是\dotfill(\kh{B})
\xx{$\tan{x}$}{$\tan{y}$}{$\tan(x+y)$}{$\tan(x-y)$}

\item 化简: $\dfrac{\tan(\frac{\pi}{4}+\alpha)-\tan(\frac{\pi}{4}-\alpha)}{\tan2\alpha}-\tan(\frac{\pi}{4}+\alpha)\tan(\frac{\pi}{4}-\alpha)$. 
\begin{jieda}
$2\alpha=\dfrac{\pi}{4}+\alpha-\left(\dfrac{\pi}{4}-\alpha\right)$, 于是只需要展开$\tan2\alpha$即可. \\
\da 1. 
\end{jieda}

\item 求证: $\tan5\alpha-\tan3\alpha-\tan2\alpha=\tan5\alpha\tan3\alpha\tan2\alpha$. 
\begin{jieda}
$\iff \tan5\alpha=\dfrac{\tan3\alpha+\tan2\alpha}{1-\tan3\alpha\tan2\alpha}$. \\
\ps 另外, 由第10题的结论立得. 
\end{jieda}


\item 已知关于$x$的方程$mx^2+(2m-3)x+(m-2)=0(m\neq 0)$的两个根为$\tan\alpha,\tan\beta$, 求$\tan(\alpha+\beta)$的最小值.
\begin{jieda}
$\Delta\geq 0\land m\neq 0\iff m\leq\dfrac{9}{4}\land m\neq 0$. \\
$\tan(\alpha+\beta)=\dfrac{\tan\alpha+\tan\beta}{1-\tan\alpha\tan\beta}=\dfrac{3}{2}-m\geq-\dfrac{3}{4}$. \\
故$\min=-\dfrac{3}{4}$, $m=\dfrac{9}{4}$. 
\end{jieda}

\item 设$A,B,C$是$\triangle{ABC}$的三内角, 已知$\dfrac{\tan{A}-\tan{B}}{\tan{A}+\tan{B}}=\dfrac{\sin{C}-\sin{B}}{\sin{C}}$, 
求$\cos\dfrac{B+C}{2}$的值. 
\begin{jieda}
切割化弦, 通分得
\begin{align*}
\dfrac{\sin{C}-\sin{B}}{\sin{C}}
& =\dfrac{\sin{A}\cos{B}-\sin{B}\cos{A}}{\sin{A}\cos{B}+\sin{B}\cos{A}}\\
& =\dfrac{\sin(A-B)}{\sin(A+B)}\\
& =\dfrac{\sin(A-B)}{\sin{C}}
\end{align*}
故$\sin{C}-\sin{B}=\sin(A-B)$, 即$\sin{B}=\sin(A+B)-\sin(A-B)=2\cos{A}\sin{B}$. \\
又$B\in(0,\pi)$, 则$\sin{B}\neq 0$, 故$\cos{A}=\dfrac{1}{2}$, 即$A=\dfrac{\pi}{3}$. \\
故$\cos\dfrac{B+C}{2}=\cos\dfrac{\pi}{3}=\dfrac{1}{2}$. 
\end{jieda}
\end{enumerate}




\subsection{两角和与差的余弦、正弦和正切公式(4)}
\zsd 综合应用


\begin{enumerate}
\item 在$\triangle{ABC}$中, 已知$\sin{A}=\dfrac{5}{13}$, $\cos{B}=\dfrac{4}{5}$, 则$\sin{C}$的值为\blkda{$\dfrac{56}{65}$}. 
\begin{jieda}
由题意可得$\cos{A}=\pm\dfrac{12}{13}$, $\sin{B}=\dfrac{3}{5}$. \\
$\sin{C}=\sin(\pi-A-B)=\sin(A+B)=\sin{A}\cos{B}+\cos{A}\sin{B}$. \\
当$\cos{A}=\dfrac{12}{13}$时, $\sin{C}=\dfrac{56}{65}$; \\
当$\cos{A}=-\dfrac{12}{13}$时, $\sin{C}=-\dfrac{16}{65}<0$, 故舍去. 
\end{jieda}


\item 已知$\alpha,\beta$为锐角, 且$\sin\alpha=\dfrac{8}{17}$, $\cos(\alpha-\beta)=\dfrac{21}{29}$, 
则$\cos\beta$的值为\blkda{$\dfrac{155}{493}$}.  % $\dfrac{415}{493}$不符合条件. 
\begin{jieda}
$\cos\alpha=\dfrac{15}{17}$, $\sin(\alpha-\beta)=\pm\dfrac{20}{29}$. \\
$\cos\beta=\cos(-\beta)=\cos(\alpha-\beta-\alpha)=\cos(\alpha-\beta)\cos\alpha+\sin(\alpha-\beta)\sin\beta$. \\
当$\sin(\alpha-\beta)=\dfrac{20}{29}$时, $\sin\beta=-\dfrac{132}{493}<0$, 故舍去; \\ % $\cos\beta=\dfrac{415}{493}$; \\
当$\sin(\alpha-\beta)=-\dfrac{20}{29}$时, $\sin\beta=\dfrac{468}{493}$, $\cos\beta=\dfrac{155}{493}$. \\
\ps $\alpha=\arcsin\dfrac{8}{17}$, $\alpha-\beta=\pm\arccos\dfrac{21}{29}$. \\
若$\alpha-\beta=\arccos\dfrac{21}{29}$, 则$\beta=\arcsin\dfrac{8}{17}-\arccos\dfrac{21}{29}$, 计算器可知$\beta<0$, 故舍去. \\
若$\alpha-\beta=-\arccos\dfrac{21}{29}$, 则$\beta=\arcsin\dfrac{8}{17}+\arccos\dfrac{21}{29}$, 计算器可知$\beta$为锐角, 符合条件. 
\end{jieda}

\item 已知: $\sin\alpha+\sin\beta=\dfrac{1}{2}$, $\cos\alpha+\cos\beta=\dfrac{1}{3}$, 则$\cos(\alpha-\beta)$=\blkda{$-\dfrac{59}{72}$}. 

\item 化简: $2\sin\left(\dfrac{\pi}{4}+\alpha\right)\sin\left(\dfrac{\pi}{4}-\alpha\right)$=\blkda{$\cos2\alpha$}. 

\item 设$A,B,C$为$\triangle{ABC}$的三内角, 则在下列表达式中, 一定是常数的是\blkda{\ding{172}\ding{174}\ding{175}}: \\
\ding{172} $\cos(A+B)+\cos{C}$;
\ding{173} $\sin(A+B)+\sin{C}$;
\ding{174} $\tan\dfrac{A+B}{2}\tan\dfrac{C}{2}$;
\ding{175} $\cot\dfrac{A+B}{2}\cot\dfrac{C}{2}+\cot\dfrac{A+C}{2}\cot\dfrac{B}{2}+\cot\dfrac{B+C}{2}\cot\dfrac{A}{2}$. 

\item 若$\alpha+\beta=\dfrac{\pi}{4}$, $\alpha,\beta\in\left(0,\dfrac{\pi}{2}\right)$, 
则$(1+\tan\alpha)(1+\tan\beta)$的值为\blkda{2}. 
\begin{jieda}
$1=\tan(\alpha+\beta)=\dfrac{\tan\alpha+\tan\beta}{1-\tan\alpha\tan\beta}$, 故$1-\tan\alpha\tan\beta=\tan\alpha+\tan\beta$, \\
于是$(1+\tan\alpha)(1+\tan\beta)=2$. \\
\ps $\alpha+\beta=\dfrac{\pi}{4}\iff (1+\tan\alpha)(1+\tan\beta)=2$, 只要正切有意义就是成立的. 
\end{jieda}


\item 已知关于$x$的方程$x^2+3ax+3a+1=0(a>2)$的两根分别是$\tan\alpha,\tan\beta$, 且$\alpha,\beta\in\left(-\dfrac{\pi}{2},\dfrac{\pi}{2}\right)$, 
则$\alpha+\beta$=\blkda{$-\dfrac{3\pi}{4}$}. 
\begin{jieda}
$\Delta\geq 0\land a>2\iff a>2$. \\
故$\begin{cases}
\tan\alpha+\tan\beta=-3a<0\\
\tan\alpha\tan\beta=3a+1>0
\end{cases}$, 
故$\tan\alpha,\tan\beta<0$. \\
又$\alpha,\beta\in\left(-\dfrac{\pi}{2},\dfrac{\pi}{2}\right)$, 故$\alpha,\beta\in\left(-\dfrac{\pi}{2},0\right)$, 则$\alpha+\beta\in(-\pi,0)$. \\
则$\tan(\alpha+\beta)=\dfrac{\tan\alpha+\tan\beta}{1-\tan\alpha\tan\beta}=\dfrac{-3a}{1-(3a+1)}=1$, 故$\alpha+\beta=-\dfrac{3\pi}{4}$. 
\end{jieda}


\item 已知$\cos\alpha\cos\beta=\dfrac{1}{5}$, 则$\sin\alpha\sin\beta$的取值范围是\blkda{$\left[-\dfrac{4}{5},\dfrac{4}{5}\right]$}. 
\begin{jieda}
令$t=\sin\alpha\sin\beta$, 则
$\begin{cases}
\frac{1}{5}+t=\cos(\alpha-\beta)\in[-1,1]\\
\frac{1}{5}-t=\cos(\alpha+\beta)\in[-1,1]
\end{cases}$, 
解得$t\in[-\frac{4}{5},\frac{4}{5}]$. \\
令$\alpha=\beta=\arccos\frac{1}{\sqrt{5}}$, 则$t=\frac{4}{5}$; 令$\alpha=\arccos\frac{1}{\sqrt{5}},\beta=-\arccos\frac{1}{\sqrt{5}}$, 则$t=-\frac{4}{5}$. 
\end{jieda}

\item 在$\triangle{ABC}$中, 若$\tan{A}\cdot\tan{B}>1$, 那么$\triangle{ABC}$是\dotfill(\kh{A})
\xx{锐角三角形}{直角三角形}{钝角三角形}{以上三种情况均有可能}

\item 已知$\theta$是第一象限角, 那么恒有\dotfill(\kh{D})
\xx{$\sin\dfrac{\theta}{2}>0$}
{$\sin\dfrac{\theta}{2}<\cos\dfrac{\theta}{2}$}
{$\sin\dfrac{\theta}{2}>\cos\dfrac{\theta}{2}$}
{$\tan\dfrac{\theta}{2}<1$}
\begin{jieda}
\begin{minipage}{\textwidth}
\centering
\includegraphics[scale=0.8]{./bodyxb/src/main42chap05renyijiaodesanjiaobi01.pdf}
\end{minipage}
结合图像可知, $\sin\dfrac{\theta}{2}$可正可负, $\sin\dfrac{\theta}{2}-\cos\dfrac{\theta}{2}$可正可负, $0<\tan\dfrac{\theta}{2}<1$. 
\end{jieda}



\item 已知$\begin{cases}
1+\cos\alpha-\sin\beta+\sin\alpha\sin\beta=0\\
1-\cos\alpha-\cos\beta+\sin\alpha\cos\beta=0
\end{cases}$, 求$\sin\alpha$的值. 
\begin{jieda}
$\begin{cases}
1+\cos\alpha=(1-\sin\alpha)\sin\beta\\
1-\cos\alpha=(1-\sin\alpha)\cos\beta
\end{cases}$, 
两式两边平方再相加可得$2+2\cos^2\alpha=(1-\sin\alpha)^2$. \\
即$3\sin^2\alpha-2\sin\alpha-3=0$, 解得$\sin\alpha=\dfrac{1-\sqrt{10}}{3}$(另一个大于1, 已舍去). 
\end{jieda}


\item 已知$\cos\left(\alpha-\dfrac{\beta}{2}\right)=\dfrac{12}{13}$, $\sin\left(\beta-\dfrac{\alpha}{2}\right)=\dfrac{4}{5}$, 
$\alpha,\beta\in\left(0,\dfrac{\pi}{2}\right)$, 求$\cos(\alpha+\beta)$的值. 
\begin{jieda}
$\alpha-\dfrac{\beta}{2},\beta-\dfrac{\alpha}{2}\in\left(-\dfrac{\pi}{4},\dfrac{\pi}{2}\right)$. \\
因为$\sin\left(\beta-\dfrac{\alpha}{2}\right)=\dfrac{4}{5}>\dfrac{\sqrt{2}}{2}$, 故$\beta-\dfrac{\alpha}{2}\in\left(\dfrac{\pi}{4},\dfrac{\pi}{2}\right)$, 
且$\cos\left(\beta-\dfrac{\alpha}{2}\right)=\dfrac{3}{5}$. \\
因为$\cos\left(\alpha-\dfrac{\beta}{2}\right)=\dfrac{12}{13}$, 故$\sin\left(\alpha-\dfrac{\beta}{2}\right)=\pm\dfrac{5}{13}$. \\
又$\dfrac{\alpha+\beta}{2}=\alpha-\dfrac{\beta}{2}+\beta-\dfrac{\alpha}{2}$, 且$\cos(\alpha+\beta)=2\cos^2\dfrac{\alpha+\beta}{2}-1$, 故
\begin{itemize}
\item 当$\sin\left(\alpha-\dfrac{\beta}{2}\right)=\dfrac{5}{13}$时, $\cos\dfrac{\alpha+\beta}{2}=\dfrac{16}{65}$, 故$\cos(\alpha+\beta)=-\dfrac{3713}{4225}$; 
\item 当$\sin\left(\alpha-\dfrac{\beta}{2}\right)=-\dfrac{5}{13}$时, $\cos\dfrac{\alpha+\beta}{2}=\dfrac{56}{65}$, 故$\cos(\alpha+\beta)=\dfrac{2047}{4225}$. 
\end{itemize}
\ps $\alpha-\dfrac{\beta}{2}=\pm\arccos\dfrac{12}{13}$, $\beta-\dfrac{\alpha}{2}=\arcsin\dfrac{4}{5}$. \\
若$\alpha-\dfrac{\beta}{2}=\arccos\dfrac{12}{13}$, 
则$\alpha=\dfrac{2}{3}\left(\arcsin\dfrac{4}{5}+2\arccos\dfrac{12}{13}\right)$, 
$\beta=\dfrac{2}{3}\left(2\arcsin\dfrac{4}{5}+\arccos\dfrac{12}{13}\right)$, 计算器可知符合条件. \\
若$\alpha-\dfrac{\beta}{2}=-\arccos\dfrac{12}{13}$, 
则$\alpha=\dfrac{2}{3}\left(\arcsin\dfrac{4}{5}-2\arccos\dfrac{12}{13}\right)$, 
$\beta=\dfrac{2}{3}\left(2\arcsin\dfrac{4}{5}-\arccos\dfrac{12}{13}\right)$, 计算器可知符合条件. 
\end{jieda}

\item 已知$\tan(\alpha-\beta)=\dfrac{1}{2}$, $\tan\beta=-\dfrac{1}{7}$, 且$\alpha,\beta\in(0,\pi)$, 求$2\alpha-\beta$的大小. 
\begin{jieda}
$\tan\beta=-\dfrac{1}{7}<-1\land\beta\in(0,\pi)$, 则$\beta\in\left(\dfrac{3\pi}{4},\pi\right)$. \\
$\tan\alpha=\tan(\alpha-\beta+\beta)=\dfrac{\frac{1}{2}-\frac{1}{7}}{1+\frac{1}{2}\cdot\frac{1}{7}}=\dfrac{1}{3}<1$, 
又$\alpha\in(0,\pi)$, 故$\alpha\in\left(0,\dfrac{\pi}{4}\right)$. \\
于是$2\alpha-\beta\in\left(-\pi,-\dfrac{\pi}{4}\right)$. \\
则$\tan(2\alpha-\beta)=\tan(\alpha-\beta+\alpha)=\dfrac{\frac{1}{2}+\frac{1}{3}}{1-\frac{1}{2}\cdot\frac{1}{3}}=1$, 
故$2\alpha-\beta=-\dfrac{3\pi}{4}$. 
\end{jieda}


\item
\begin{enumerate}[(1)]
\item 已知$\triangle{ABC}$不是直角三角形, 求证: $\tan{nA}+\tan{nB}+\tan{nC}=\tan{nA}\tan{nB}\tan{nC}(n\zz)$; 
\item 在$\triangle{ABC}$中, 求证: $\tan\dfrac{A}{2}\tan\dfrac{B}{2}+\tan\dfrac{B}{2}\tan\dfrac{C}{2}+\tan\dfrac{C}{2}\tan\dfrac{A}{2}=1$. 
\end{enumerate}
\begin{jieda}
\begin{enumerate}[(1)]
\item $n(A+B+C)=n\pi$, 故$nC=n\pi-nA-nB$. \\
于是$\tan{nC}=\tan(n\pi-nA-nB)=-\tan(nA+nB)=-\dfrac{\tan{nA}+\tan{nB}}{1-\tan{nA}\tan{nB}}$, \\
则$\tan{nC}(1-\tan{nA}\tan{nB})=-\tan{nA}-\tan{nB}$, 整理即得. 

\item $\iff\tan\dfrac{C}{2}\cdot\left(\tan\dfrac{A}{2}+\tan\dfrac{B}{2}\right)=1-\tan\dfrac{A}{2}\tan\dfrac{B}{2}
\iff\tan\dfrac{C}{2}=\dfrac{1-\tan\frac{A}{2}\tan\frac{B}{2}}{\tan\frac{A}{2}+\tan\frac{B}{2}}$. \\
$\LHS=\tan\dfrac{\pi-A-B}{2}=\cot\dfrac{A+B}{2}=\dfrac{1}{\tan\frac{A+B}{2}}=\RHS$. 
\end{enumerate}
\end{jieda}



\end{enumerate}




\subsection{两角和与差的余弦、正弦和正切公式(5)}
\zsd 辅助角公式

用辅助角公式化简
\begin{enumerate}
\item $\dfrac{\sqrt{3}}{2}\sin\alpha+\dfrac{1}{2}\cos\alpha$=\blkda[3]{$\sin\left(\alpha+\dfrac{\pi}{6}\right)$}. 

\item $\dfrac{\sqrt{3}}{2}\sin\alpha-\dfrac{1}{2}\cos\alpha$=\blkda[3]{$\sin\left(\alpha-\dfrac{\pi}{6}\right)$}. 

\item $\sin\alpha-\sqrt{3}\cos\alpha$=\blkda[3]{$2\sin\left(\alpha-\dfrac{\pi}{3}\right)$}. 

\item $2\cos\alpha+2\sqrt{3}\sin\alpha$=\blkda[3]{$4\sin\left(\alpha+\dfrac{\pi}{6}\right)$}. 

\item $\sin\alpha+\cos\alpha$=\blkda[3]{$\sqrt{2}\sin\left(\alpha+\dfrac{\pi}{4}\right)$}. 

\item $\sin\alpha-\cos\alpha$=\blkda[3]{$\sqrt{2}\sin\left(\alpha-\dfrac{\pi}{4}\right)$}. 

\item $-\sin\alpha-\cos\alpha$=\blkda[3]{$\sqrt{2}\sin\left(\alpha+\dfrac{5\pi}{4}\right)$}. 

\item $3\sin\alpha-4\cos\alpha$=\blkda[4]{$5\sin\left(\alpha-\arccos\dfrac{3}{5}\right)$}. 
\begin{jieda}
$\ang{53}$不可以, 因为$\ang{53}\neq\arccos\frac{3}{5}$. 
\end{jieda}

\item $\sqrt{3}\sin\left(\alpha+\dfrac{\pi}{6}\right)-3\cos\left(\alpha+\dfrac{\pi}{6}\right)$=\blkda[3]{$2\sqrt{3}\sin\left(\alpha-\dfrac{\pi}{6}\right)$}. 

\item $\dfrac{\sqrt{2}}{6}\sin\left(\dfrac{\pi}{3}-\alpha\right)+\dfrac{\sqrt{6}}{6}\cos\left(\dfrac{\pi}{3}-\alpha\right)$
=\blkda[3]{$\dfrac{\sqrt{2}}{3}\sin\left(\alpha+\dfrac{\pi}{3}\right)$}. 


\item 将下列各式化为$A\sin(\alpha+\varphi)(A>0)$的形式: 
\begin{enumerate}[(1)]
\item $\cos\left(\dfrac{\pi}{3}+x\right)+\cos{x}$; 

\item $\sin\left(\dfrac{\pi}{3}-2x\right)+\cos\left(\dfrac{\pi}{3}+2x\right)$. 
\end{enumerate}
\begin{jieda}
\begin{enumerate}[(1)]
\item $E=\cos\dfrac{\pi}{3}\cos{x}-\sin\dfrac{\pi}{3}\sin{x}+\cos{x}=\dfrac{3}{2}\cos{x}-\dfrac{\sqrt{3}}{2}\sin{x}=\sqrt{3}\sin\left(x+\dfrac{2\pi}{3}\right)$. 

\item $E=\sin\dfrac{\pi}{3}\cos2x-\cos\dfrac{\pi}{3}\sin2x+\cos\dfrac{\pi}{3}\cos2x-\sin\dfrac{\pi}{3}\sin2x=\dfrac{\sqrt{6}+\sqrt{2}}{2}\sin\left(2x+\dfrac{3\pi}{4}\right)$. 
\end{enumerate}
\end{jieda}

\item 已知$m=\dfrac{\sqrt{5}\sin\alpha+1}{\cos\alpha+2}$, 求$m$的取值范围. 
\begin{jieda}
$m\cos\alpha+2m=\sqrt{5}\sin\alpha+1$, 即$\sqrt{5}\sin\alpha-m\cos\alpha=2m-1$, \\
故$\sqrt{5+m^2}\sin(\alpha-\varphi)=2m-1$. \\
于是$-1\leq \dfrac{2m-1}{\sqrt{5+m^2}}\leq 1$, 解得$m\in\left[-\dfrac{2}{3},2\right]$. 
\end{jieda}

\item 已知$n=3\sin(x+\ang{20})+5\sin(x+\ang{80})$, 求$n$的最大值. 
\begin{jieda}
令$t=x+\ang{20}$, 则
\begin{align*}
n & =3\sin{t}+5\sin(t+\ang{60})\\
& = 3\sin{t}+5\sin{t}\cos\ang{60}+5\cos{t}\sin\ang{60}\\
& = \dfrac{11}{2}\sin{t}+\dfrac{5\sqrt{3}}{2}\cos{t}\\
& = 7\sin\left(t+\arccos\dfrac{11}{14}\right)
\end{align*}
故$n_{\max}=7$. 
\end{jieda}

\item 已知$5\sqrt{3}\sin{x}+5\cos{x}=6$, $\sqrt{2}\sin{y}+\sqrt{6}\cos{y}=1$, 
且$x\in\left(0,\dfrac{\pi}{3}\right)$, $y\in\left(\dfrac{\pi}{6},\dfrac{\pi}{2}\right)$, 
求$\cos(x+y)$. 
\begin{jieda}
$\iff\begin{cases}
\sin(x+\frac{\pi}{6})=\frac{3}{5}\\
\sin(y+\frac{\pi}{3})=\frac{1}{2\sqrt{2}}
\end{cases}$. \\
又$x+\frac{\pi}{6}\in(\frac{\pi}{6},\frac{\pi}{2})$, $y+\frac{\pi}{3}\in(\frac{\pi}{2},\frac{5\pi}{6})$, 
故$\begin{cases}
\cos(x+\frac{\pi}{6})=\frac{4}{5}\\
\cos(y+\frac{\pi}{3})=-\frac{\sqrt{7}}{2\sqrt{2}}
\end{cases}$. \\
故$E=\cos(x+\frac{\pi}{6}+y+\frac{\pi}{3}-\frac{\pi}{2})=\sin(x+\frac{\pi}{6}+y+\frac{\pi}{3})=\dfrac{4\sqrt{2}-3\sqrt{14}}{20}$. 
\end{jieda}

\end{enumerate}





\subsection{二倍角公式}
\zsd 二倍角正弦、余弦和正切公式


\begin{enumerate}
\item $\sqrt{2}\sin\dfrac{\pi}{32}\cos\dfrac{\pi}{32}\cos\dfrac{\pi}{16}\cos\dfrac{\pi}{8}$=\blkda{$\dfrac{1}{8}$}. 

\item 已知$\cos(\pi+\alpha)=\dfrac{1}{3}(\pi<\alpha<2\pi)$, 则$\sin2\alpha$=\blkda{$\dfrac{4\sqrt{2}}{9}$}. 
\begin{jieda}
$\cos\alpha=-\dfrac{1}{3}$, $\pi<\alpha<2\pi$, 故$\alpha$是第三象限角. \\
法一: $\sin2\alpha=2\sin\alpha\cos\alpha=2\cdot\dfrac{-2\sqrt{2}}{3}\cdot\dfrac{-1}{3}$. \\
法二: $\alpha=2\pi-\arccos(-\frac{1}{3})=\pi+\arccos\frac{1}{3}$, 然后计算器可以解决. 
\end{jieda}


\item 若$\alpha\in(0,\pi)$, $\sin\alpha+\cos\alpha=\dfrac{1}{3}$, 则$\cos2\alpha$=\blkda{$-\dfrac{\sqrt{17}}{9}$}. 
\begin{jieda}
法一: $\dfrac{1}{9}=(\sin\alpha+\cos\alpha)^2=1+2\sin\alpha\cos\alpha$, 则$2\sin\alpha\cos\alpha=-\dfrac{8}{9}<0$. \\
又$\alpha\in(0,\pi)$, 故$\sin\alpha>0>\cos\alpha$. \\
于是$\cos\alpha-\sin\alpha=-\sqrt{(\cos\alpha-\sin\alpha)^2}=-\sqrt{1-2\cos\alpha\sin\alpha}=-\dfrac{\sqrt{17}}{3}$. \\
故$\cos^2\alpha=(\cos\alpha-\sin\alpha)(\cos\alpha+\sin\alpha)=-\dfrac{\sqrt{17}}{9}$. \\
法二: $\begin{cases}
\sin\alpha+\cos\alpha=\frac{1}{3}\\
\sin^2\alpha+\cos^2\alpha=1
\end{cases}$且$\alpha\in(0,\pi)$, 解得$\cos\alpha=\frac{1-\sqrt{17}}{6},\sin\alpha=\frac{1+\sqrt{17}}{6}$. \\
法三: $\sqrt{2}\sin(\alpha+\frac{\pi}{4})=\frac{1}{3}$, 且$\alpha\in(0,\pi)$, 
故$\alpha+\frac{\pi}{4}=\pi-\arcsin\frac{1}{3\sqrt{2}}$($\arcsin\frac{1}{3\sqrt{2}}$舍去).\\
于是$\alpha=\frac{3\pi}{4}-\arcsin\frac{1}{3\sqrt{2}}$. 
\end{jieda}

\item 若$\theta\in\left(\dfrac{\pi}{4},\dfrac{\pi}{2}\right)$, $\sin2\theta=\dfrac{1}{16}$, 则$\cos\theta-\sin\theta$ =\blkda{$-\dfrac{\sqrt{15}}{4}$}. 
\begin{jieda}
因为$\theta\in\left(\dfrac{\pi}{4},\dfrac{\pi}{2}\right)$, 故$\cos\theta<\sin\theta$. \\
于是$\cos\theta-\sin\theta=-\sqrt{(\cos\theta-\sin\theta)^2}=-\sqrt{1-\sin2\theta}=-\dfrac{\sqrt{15}}{4}$. 
\end{jieda}

\item 已知角$\theta$的顶点与坐标原点重合, 始边与$x$轴的正半轴重合, 终边在射线$y=2x(x\geq 0)$上, 则$\cos2\theta$=\blkda{$-\dfrac{3}{5}$}. 
\begin{jieda}
$\tan\theta=2$, 则$\cos2\theta=\dfrac{1-\tan^2\theta}{1+\tan^2\theta}=-\dfrac{3}{5}$
\end{jieda}


\item 已知$\dfrac{3\pi}{2}<\theta<2\pi$, 化简: $\sqrt{\dfrac{1}{2}-\dfrac{1}{2}\sqrt{\dfrac{1}{2}+\dfrac{1}{2}\cos2\theta}}$ =\blkda{$\sin\dfrac{\theta}{2}$}. 

\item 若$\dfrac{\sin\alpha+\cos\alpha}{\sin\alpha-\cos\alpha}=\dfrac{1}{2}$, 则$\tan2\alpha$=\blkda{$\dfrac{3}{4}$}. 
\begin{jieda}
齐次化后得$\tan\alpha=-3$, 故$\tan2\alpha=\dfrac{2\tan\alpha}{1-\tan^2\alpha}=\dfrac{3}{4}$. 
\end{jieda}

\item 已知$\sin\alpha\cos\alpha=-\dfrac{1}{8}$, 且$\dfrac{\pi}{2}<\alpha<\dfrac{3\pi}{2}$, 则$\sin\alpha-\cos\alpha$=\blkda{$\dfrac{\sqrt{5}}{2}$}. 
\begin{jieda}
因为$\sin\alpha\cos\alpha<0$, 且$\dfrac{\pi}{2}<\alpha<\dfrac{3\pi}{2}$, 故$\alpha\in\left(\dfrac{\pi}{2},\pi\right)$, 则$\sin\alpha>0>\cos\alpha$. 
\end{jieda}

\item 若$\sin\left(\dfrac{\pi}{4}-x\right)=\dfrac{5}{13}$, 且$0<x<\dfrac{\pi}{4}$, 则$\dfrac{\cos2x}{\cos(\frac{\pi}{4}+x)}=$\dotfill(\kh{C})
\xx{$\dfrac{13}{24}$}{$\dfrac{12}{13}$}{$\dfrac{24}{13}$}{$\dfrac{13}{12}$}

% 1978高考
\item 已知$\alpha,\beta$都是锐角, 且$3\sin^2\alpha+2\sin^2\beta=1$, $3\sin2\alpha-2\sin2\beta=0$, 那么$\alpha,\beta$之间的关系是\dotfill(\kh{D})
\xx{$\alpha+\beta=\dfrac{\pi}{4}$}{$\alpha-\beta=\dfrac{\pi}{4}$}{$\alpha+2\beta=\dfrac{\pi}{4}$}{$\alpha+2\beta=\dfrac{\pi}{2}$}
\begin{jieda}
法一: $3\sin^2\alpha=1-2\sin^2\beta=\cos2\beta$, $3\sin\alpha\cos\alpha=\sin2\beta$, 故$\tan2\beta=\cot\alpha$. \\
法二: $3\sin^2\alpha=1-2\sin^2\beta=\cos2\beta$, $3\sin\alpha\cos\alpha=\sin2\beta$,\\
故$\left(3\sin^2\alpha\right)^2+\left(3\sin\alpha\cos\alpha\right)^2=\cos^22\beta+\sin^22\beta=1$, \\
解得$\sin\alpha=\dfrac{1}{3}=\cos2\beta$. 
\end{jieda}

\item 设$\sin\theta+\cos\theta=\dfrac{\sqrt{2}}{3}$, $\dfrac{\pi}{2}<\theta<\pi$, 求$\tan\theta-\cot\theta$的值. 
\begin{jieda}
由条件可得$\sin\theta\cos\theta=-\dfrac{7}{18}$, $\sin\theta-\cos\theta=\dfrac{4}{3}$. 
\begin{align*}
E & = \dfrac{\sin\theta}{\cos\theta}-\dfrac{\cos\theta}{\sin\theta}\\
& = \dfrac{\sin^2\theta-\cos^2\theta}{\sin\theta\cos\theta}\\
& = \dfrac{(\sin\theta-\cos\theta)(\sin\theta+\cos\theta)}{\sin\theta\cos\theta}\\
& = -\dfrac{8\sqrt{2}}{7}. 
\end{align*}
\ps $2\sin\theta\cos\theta=(\sin\theta+\cos\theta)^2-1=1-(\sin\theta-\cos\theta)^2$. 
\end{jieda}


\item 已知$\tan\left(\dfrac{\pi}{4}+\alpha\right)=\dfrac{1}{2}$. 
\begin{enumerate}[(1)]
\item 求$\tan\alpha$的值; 
\item 求$\dfrac{\sin2\alpha-\cos^2\alpha}{1+\cos2\alpha}$的值. 
\end{enumerate}
\begin{jieda}
\begin{enumerate}[(1)]
\item $\tan\alpha=\tan\left(\alpha+\dfrac{\pi}{4}-\dfrac{\pi}{4}\right)=-\dfrac{1}{3}$. 

\item 齐次化: $E=\dfrac{2\sin\alpha\cos\alpha-\cos^2\alpha}{1+\cos^2\alpha-\sin^2\alpha}=\tan\alpha-\dfrac{1}{2}=-\dfrac{5}{6}$. 
\end{enumerate}
\end{jieda}


\item 已知$6\sin^2\alpha+\sin\alpha\cos\alpha-2\cos^2\alpha=0$, $\alpha\in\left[\dfrac{\pi}{2},\pi\right]$, 求$\sin\left(2\alpha+\dfrac{\pi}{3}\right)$的值. 
\begin{jieda}
$\iff 6\tan^2\alpha+\tan\alpha-2=0$. 又$\alpha\in\left[\dfrac{\pi}{2},\pi\right]$, 故$\tan\alpha=-\dfrac{2}{3}$. \\
故$\sin\left(2\alpha+\dfrac{\pi}{3}\right)=\sin2\alpha\cos\dfrac{\pi}{3}+\cos2\alpha\sin\dfrac{\pi}{3}=f(\tan\alpha)=\dfrac{5\sqrt{3}-12}{26}$. 
\end{jieda}

\item 对于集合$S=\{\theta_1,\theta_2,\cdots,\theta_n\}$和常数$\theta_0$, 
定义$\mu=\dfrac{\cos^2(\theta_1-\theta_0)+\cos^2(\theta_2-\theta_0)+\cdots+\cos^2(\theta_n-\theta_0)}{n}$为集合$S$相对于$\theta_0$的``余弦方差'', 
求证: 集合$\left\{\dfrac{\pi}{3},\dfrac{2\pi}{3},\pi\right\}$相对于任何常数$\alpha$的``余弦方差''是一个与$\alpha$无关的常数. 
\begin{jieda}
等价于证明$\cos^2\left(\dfrac{\pi}{3}-\alpha\right)+\cos^2\left(\dfrac{2\pi}{3}-\alpha\right)+\cos^2(\pi-\alpha)$是与$\alpha$无关的常数. \\
降次, 展开, 化简, 立得所求为$\dfrac{1}{2}$. 
\end{jieda}


\item (选做)已知$f(\theta)=\sin^2\theta+\sin^2(\theta+\alpha)+\sin^2(\theta+\beta)$, 
其中$\alpha,\beta$是满足$0<\alpha<\beta<\pi$的常数, 试问: 是否存在这样的$\alpha,\beta$, 使得$f(\theta)$的值是与$\theta$无关的定值? 
若存在, 求出所有的$\alpha$和$\beta$; 若不存在, 请说明理由. 
\begin{jieda}
法一: 
$\begin{cases}
f(0)=\sin^2\alpha+\sin^2\beta\\
f(\frac{\pi}{2})=1+\cos^2\alpha+\cos^2\beta
\end{cases}$, 于是$f(0)+f(\frac{\pi}{2})=3$, 故$f(0)=f(\frac{\pi}{2})=\frac{3}{2}$. \\
又$f(-\alpha)=\sin^2\alpha+\sin^2(\beta-\alpha)$, $f(-\beta)=\sin^2\beta+\sin^2(\alpha-\beta)$, 
故$\sin^2\beta=\sin^2\alpha=\dfrac{3}{4}$. \\
结合$0<\alpha<\beta<\pi$可得$\alpha=\dfrac{\pi}{3}$, $\beta=\dfrac{2\pi}{3}$. \\
代入检验可知成立. \\
法二: $f(\theta)=\frac{3}{2}+\frac{1}{2}((\sin2\alpha+\sin2\beta)\sin2\theta-(1+\cos2\alpha+\cos2\beta)\cos2\theta)$. \\
于是$\begin{cases}
\sin2\alpha+\sin2\beta=0\\
1+\cos2\alpha+\cos2\beta=0
\end{cases}$. 
\end{jieda}
\end{enumerate}




\subsection{三角变换的应用(1)}
\zsd 半角公式、万能公式

\begin{tcolorbox}
\begin{itemize}
\item $\tan\theta+\cot\theta=\dfrac{2}{\sin2\theta}$, $\tan\theta-\cot\theta=-2\cot2\theta$. 
\item 角的范围条件与三角比值结合时, 要表述出来. 
\end{itemize}
\end{tcolorbox}

\begin{enumerate}
\item 在$\triangle{ABC}$中, 已知$\cos{A}=-\dfrac{3}{5}$, 则$\sin\dfrac{A}{2}$=\blkda{$\dfrac{2\sqrt{5}}{5}$}. 

\item 已知$\tan\alpha=\dfrac{4}{3}$, 且$\pi<\alpha<\dfrac{3\pi}{2}$, 则$\cos\dfrac{\alpha}{2}$=\blkda{$-\dfrac{\sqrt{5}}{5}$}. 

\item 已知$\sin\theta=-\dfrac{3}{5}$, $3\pi<\theta<\dfrac{7\pi}{2}$, 则$\tan\dfrac{\theta}{2}$=\blkda{$-3$}. 

\item 化简: $\dfrac{\sin2x}{1+\cos2x}\cdot\dfrac{\cos{x}}{1+\cos{x}}\cdot\dfrac{\sin{x}}{1-\cos{x}}$=\blkda{$1$}. 

\item 已知$25\sin^2\theta+\sin\theta-24=0$, $\theta$为第二象限角, 则$\cos\dfrac{\theta}{2}$=\blkda{$\pm\dfrac{3}{5}$}. 
\begin{jieda}
虽然$\sin\theta$只有一解, 且$\theta$的终边唯一确定, 但是$\frac{\theta}{2}$的终边有两种情形. 
\end{jieda}

\item 若$\dfrac{2\sin\theta+\cos\theta}{\sin\theta-3\cos\theta}=-5$, 则$3\cos2\theta+4\sin2\theta$=\blkda{$\dfrac{7}{5}$}. 
\begin{jieda}
$-5=\dfrac{2\tan\theta+1}{\tan\theta-3}$, 故$\tan\theta=2$. \\
于是$E=3\cdot\dfrac{1-\tan^2\theta}{1+\tan^2\theta}+4\cdot\dfrac{2\tan\theta}{1+\tan^2\theta}=\dfrac{7}{5}$. 
\end{jieda}

\item 不用计算器计算: $\tan\ang{15}+\cot\ang{15}=$\dotfill(\kh{C})
\xx{2}{$2+\sqrt3$}{4}{$\dfrac{4\sqrt{3}}{3}$}

\item 已知$2\sin{x}=1+\cos{x}$, 则$\tan\dfrac{x}{2}$的值为\dotfill(\kh{B})
\xx{$\dfrac{1}{2}$}{$\dfrac{1}{2}$或不存在}{2或$\dfrac{1}{2}$}{不存在}
\begin{jieda}
$\begin{cases}
2\sin{x}=1+\cos{x}\\
\sin^2x+\cos^2x=1
\end{cases}$, 
故$(\cos{x},\sin{x})=(-1,0),(\frac{3}{5},\frac{4}{5})$. \\
\ps 有(由)方程组, 解得. 
\end{jieda}


\item 已知$\cos(\pi-A)=\dfrac{12}{13}$, $A$是第二象限的角, 则$\cot\dfrac{3\pi-A}{2}=$\dotfill(\kh{D})
\xx{$\dfrac{1}{5}$}{$-\dfrac{1}{5}$}{$-5$}{5}
\begin{jieda}
$\cot\dfrac{3\pi-A}{2}=\tan\dfrac{A}{2}=\dfrac{1+\cos{A}}{\sin{A}}$
\end{jieda}

\item 已知$\dfrac{\tan\alpha+1}{1-\tan\alpha}=2007$, 则$\dfrac{1+\sin2\alpha}{\cos2\alpha}$的值为\dotfill(\kh{B})
\xx{2006}{2007}{2008}{2010}
\begin{jieda}
$E=\dfrac{(\cos\alpha+\sin\alpha)^2}{\cos^2\alpha-\sin^2\alpha}=\dfrac{\cos\alpha+\sin\alpha}{\cos\alpha-\sin\alpha}=\dfrac{1+\tan\alpha}{1-\tan\alpha}$. 
\end{jieda}

\item 在$\triangle{ABC}$中, $\tan\dfrac{A}{2}+\cot\dfrac{A}{2}=\dfrac{10}{3}$, $\cos{B}=\dfrac{5}{13}$. 
\begin{enumerate}[(1)]
\item 求$\cos(A-B)$的值; 
\item 求$\cos\dfrac{A-B}{2}$的值. 
\end{enumerate}
\begin{jieda}
$\dfrac{10}{3}=\tan\dfrac{A}{2}+\cot\dfrac{A}{2}=\dfrac{2}{\sin{A}}$, 故$\sin{A}=\dfrac{3}{5}$. \\
又$\sin{B}=\dfrac{12}{13}>\sin{A}$, 于是$0<A<B<\dfrac{\pi}{2}$. \\
\begin{enumerate*}[(1)]
\item $\cos(A-B)=\dfrac{56}{65}$;  
\item $\cos\dfrac{A-B}{2}=\dfrac{11\sqrt{130}}{130}$. 
\end{enumerate*}
\end{jieda}

\item 已知$\sin\left(\dfrac{\pi}{4}+2\alpha\right)\cdot\sin\left(\dfrac{\pi}{4}-2\alpha\right)=\dfrac{1}{4}$, $\alpha\in\left(\dfrac{\pi}{4},\dfrac{\pi}{2}\right)$, 
求$2\sin^2\alpha+\tan\alpha-\cot\alpha-1$的值. 
\begin{jieda}
法一: 因为$\sin(x+y)\sin(x-y)=\sin^2x-\sin^2y$, 
故$\sin^2\dfrac{\pi}{4}-\sin^22\alpha=\dfrac{1}{4}$, 即$\sin^22\alpha=\dfrac{1}{4}$. \\
又$\alpha\in\left(\dfrac{\pi}{4},\dfrac{\pi}{2}\right)$, 于是$\sin2\alpha=\dfrac{1}{2}$, 故$2\alpha=\dfrac{5\pi}{6}$. \\
故$E=-\cos2\alpha-2\cot2\alpha=\dfrac{5\sqrt{3}}{2}$. \\
法二: 因为$\dfrac{\pi}{4}+2\alpha+\dfrac{\pi}{4}-2\alpha=\dfrac{\pi}{2}$, \\
故$2\sin\left(\dfrac{\pi}{4}+2\alpha\right)\cdot\sin\left(\dfrac{\pi}{4}-2\alpha\right)=\sin\left(2\left(\dfrac{\pi}{4}+2\alpha\right)\right)=\cos4\alpha$. \\
于是$\cos4\alpha=\dfrac{1}{2}$, 且$4\alpha\in\left(\pi,2\pi\right)$, 故$4\alpha=\dfrac{5\pi}{3}$, 即$\alpha=\dfrac{5\pi}{12}$. 
\end{jieda}

\item 已知$a\sin{x}+b\cos{x}=0$, $A\sin2x+B\cos2x=C$(其中$a,b$不同时为零), \\
求证: $2abA+(b^2-a^2)B+(a^2+b^2)C=0$. 
\begin{jieda}
法一: $\sin2x=2\sin{x}\cos{x}$, $\cos2x=2\cos^2{x}-1=1-2\sin^2x$. 
\begin{enumerate}[(1)]
\item 若$a=0$, 则$b\neq 0$, 故$\cos{x}=0$, 于是$\sin2x=0,\cos2x=-1$. \\
代入第二式可得$A\cdot 0+B\cdot(-1)=C$, 即$-B=C$. \\
于是$2abA+(b^2-a^2)B+(a^2+b^2)C=2\cdot0\cdot b\cdot A+(b^2-0^2)\cdot B+(0^2+b^2)\cdot C=0$. 

\item 若$b=0$, 则$a\neq 0$, 故$\sin{x}=0$, 则$\sin2x=0$, $\cos2x=1$. \\
代入第二式可得$A\cdot0+B\cdot1=C$, 即$B=C$. \\
于是$2abA+(b^2-a^2)B+(a^2+b^2)C=2\cdot a\cdot0\cdot A+(0^2-a^2)\cdot B+(a^2+0^2)\cdot C=0$. 

\item 若$ab\neq 0$, 则$\tan{x}=-\dfrac{b}{a}$, 于是由万能公式有$\sin2x=\dfrac{-2ab}{a^2+b^2}$, $\cos2x=\dfrac{a^2-b^2}{a^2+b^2}$. \\
代入第二式可得$A\cdot\dfrac{-2ab}{a^2+b^2}+B\cdot\dfrac{a^2-b^2}{a^2+b^2}=C$, 化简即得所求证. 
\end{enumerate}
综上, $2abA+(b^2-a^2)B+(a^2+b^2)C=0$. 

法二: $\sin{x}=\dfrac{b}{\sqrt{a^2+b^2}}$, $\cos{x}=-\dfrac{a}{\sqrt{a^2+b^2}}$. 即使如此, 还得分情况讨论得到前述结论. \\
代入第二个等式中整理即可得到所求证. 
\end{jieda}

\item 已知$-\dfrac{\pi}{4}<\alpha<\dfrac{\pi}{4}$, $\dfrac{\pi}{4}<\beta<\dfrac{3\pi}{4}$, 
且$\sin\left(\dfrac{3\pi}{4}+\alpha\right)=\dfrac{5}{13}$, $\cos\left(\dfrac{\pi}{4}-\beta\right)=\dfrac{3}{5}$, 
求$\sin\dfrac{\alpha-\beta}{2}$的值. 
\begin{jieda}
法一: $\cos(\alpha-\beta)=\cos\left(\dfrac{3\pi}{4}+\alpha+\dfrac{\pi}{4}-\beta-\pi\right)=\dfrac{16}{65}$. \\
$\dfrac{\pi}{2}<\alpha+\dfrac{3\pi}{4}<\pi$, $-\dfrac{\pi}{2}<\dfrac{\pi}{4}-\beta<0$, 故$\alpha-\beta\in(-\pi,0)$, \\
于是$\sin\dfrac{\alpha-\beta}{2}=-\dfrac{7\sqrt{130}}{130}$. \\
法二: 因为$\dfrac{\pi}{2}<\alpha+\dfrac{3\pi}{4}<\pi$, $-\dfrac{\pi}{2}<\dfrac{\pi}{4}-\beta<0$, \\
故$\alpha+\dfrac{3\pi}{4}=\pi-\arcsin\dfrac{5}{13}$, $\dfrac{\pi}{4}-\beta=-\arccos\dfrac{3}{5}$. \\
于是$\alpha-\beta=\dfrac{\pi}{2}-\arcsin\dfrac{5}{13}-\arccos\dfrac{3}{5}$. 剩下的是计算器. 
\end{jieda}

\item 上海某地区想设计建造一个自然保护区, 如图所示, 在一块矩形绿地$ABCD$中(其中$AB$为30米, $AD$为15米), 过道$EF$将其分为两个主要区域($E,F$分别是$AB,CD$边上动点), 
监测区为以$D$为圆心、$AD$为半径的四分之一圆, 古树区为四边形$BEFC$, 且$EF$与圆弧相切, 记切点为$G$, 且点在$ABCD$内部. \\
\begin{minipage}{\textwidth}
\centering
\includegraphics[scale=1]{./bodyxb/src/main42chap22banjiaogongshi01.png}
\end{minipage}
\begin{enumerate}[(1)]
\item 若$\angle{ADE}=20^{\circ}$, 求$EF$的长(结果精确到0.1); 
\item $E$点在线段$AB$上何处时, 才能使古树区的面积最大, 并求出最大值(结果精确到0.01). 
\end{enumerate}
\begin{jieda}
设$\angle{ADE}=\theta\in(0,\frac{\pi}{4}]$, 由条件可知$\angle{ADE}=\angle{GDE}$, 故$\angle{CDG}=\frac{\pi}{2}-2\theta$. \\
在$\triangle{FDG}$中, $\angle{DGF}=\frac{\pi}{2}$, 故$GF=15\tan(\frac{\pi}{2}-2\theta)=15\cot2\theta$. 同理$GE=15\tan\theta$. 
\begin{enumerate}
\item 由条件可得
\begin{align*}
EF & =EG+GF \\
& =15\tan\theta+15\cot2\theta\\
& =15\left(\tan\theta+\dfrac{1-\tan^2\theta}{2\tan\theta}\right)\\
& =\dfrac{15}{2}\left(\tan\theta+\dfrac{1}{\tan\theta}\right)\\
& \approx 23.3
\end{align*}

\item $AE=GE=15\tan\theta$, $DF=\dfrac{15}{\cos(\frac{\pi}{2}-2\theta)}=\dfrac{15}{\sin2\theta}$. \\
于是$S_{AEFD}=\dfrac{1}{2}\cdot\left(15\tan\theta+\dfrac{15}{\sin2\theta}\right)\cdot15
=\dfrac{225}{4}\left(3\tan\theta+\dfrac{1}{\tan\theta}\right)\geq\dfrac{225\sqrt{3}}{2}$. \\
当$\tan\theta=\frac{1}{\sqrt{3}}$, 即$\theta=\ang{30}$时等号成立, 此时$AE=5\sqrt{3}$, 古树区面积的最大值为$255.14$平方米. 
\end{enumerate}
\end{jieda}


\end{enumerate}




\subsection{三角变换的应用(2)}
\zsd 积化和差与和差化积公式

\begin{enumerate}
\item 已知$\alpha-\beta=\dfrac{\pi}{3}$, 则$\sin\alpha\sin\beta$的最大值为\blkda{$\dfrac{3}{4}$}. 

\item 已知$\cos3\theta=\dfrac{1}{4}$, 则$12\cos\theta\cos\left(\dfrac{\pi}{3}+\theta\right)\cos\left(\dfrac{\pi}{3}-\theta\right)$=\blkda{$\dfrac{3}{4}$}. 
\begin{jieda}
$\cos3\theta=4\cos^3\theta-3\cos\theta$, $\sin3\theta=3\sin\theta-4\sin^3\theta$\\
$\cos3\theta=4\cos\theta\cos\left(\dfrac{\pi}{3}+\theta\right)\cos\left(\dfrac{\pi}{3}-\theta\right)$\\
$\sin3\theta=4\sin\theta\sin\left(\dfrac{\pi}{3}+\theta\right)\sin\left(\dfrac{\pi}{3}-\theta\right)$\\
$\tan\theta=\tan\theta\tan\left(\dfrac{\pi}{3}+\theta\right)\tan\left(\dfrac{\pi}{3}-\theta\right)=\dfrac{3\tan\theta-\tan^3\theta}{1-3\tan^2\theta}$
\end{jieda}

\item 已知$\alpha,\beta$和$\alpha+\beta$均为锐角, $x=\sin(\alpha+\beta)$, $y=\sin\alpha+\sin\beta$, $z=\cos\alpha+\cos\beta$, 
则$x,y,z$的大小关系是\blkda[2]{$x<y<z$}. 
\begin{jieda}
$x=\sin\alpha\cos\beta+\cos\alpha\sin\beta$, 则$x<\sin\alpha\cdot 1+\sin\beta\cdot 1$, $x<\cos\alpha\cdot 1+\cos\beta\cdot 1$. \\
$y=2\sin\dfrac{\alpha+\beta}{2}\cos\dfrac{\alpha-\beta}{2}$, $z=2\cos\dfrac{\alpha+\beta}{2}\cos\dfrac{\alpha-\beta}{2}$. 
\end{jieda}


\item 已知$\sin\alpha\cos\beta=\dfrac{1}{2}$, 则$\cos\alpha\cos\beta$的取值范围是\blkda[2]{$\left[-\frac{\sqrt{3}}{2},\frac{\sqrt{3}}{2}\right]$}. 
\begin{jieda}
$\iff 1=\sin(\alpha+\beta)+\sin(\alpha-\beta)$. \\
令$t=\cos\alpha\cos\beta$, 则$2t=\cos(\alpha+\beta)+\cos(\alpha-\beta)$. \\
于是$1^2+(2t)^2=\cdots=2+2\cos2\beta$, 即$4t^2=1+2\cos2\beta\in[-1,3]$, 故$t\in[-\frac{\sqrt{3}}{2},\frac{\sqrt{3}}{2}]$. 
\end{jieda}

\item 若$\triangle{ABC}$是钝角三角形, 则$\cos^2A+\cos^2B+\cos^2C$的取值范围是\blkda{$(1,3)$}. 
\begin{jieda}
$E=1-2\cos{A}\cos{B}\cos{C}=1-(\cos(A+B)+\cos(A-B))\cos{C}=\cos^2C-\cos(A-B)\cos{C}+1$. \\
设$C$为钝角, 则$\cos(A-B)\in(0,1]$, $\cos{C}\in(-1,0)$. \\
于是$E\in(\cos^2C+1,\cos^2C-\cos{C}+1]$. \\
而$1+\cos^2C\in(1,2)$, $\cos^2C-\cos{C}+1\in(1,3)$. \\
当$A=B\to\dfrac{\pi}{4}$, $C\to\dfrac{\pi}{2}$时, $E\to 1$; \\
当$A=B\to0$, $C\to\pi$时, $E\to 3$. \\
故$E\in(1,3)$. 
\end{jieda}


\item 已知$\dfrac{1}{\sin(\ang{120}-x)}+\dfrac{1}{\sin(\ang{120}+x)}=\dfrac{4\sqrt{3}}{3}$, 则$\cos{x}$=\blkda{1或$-\dfrac{1}{4}$}. 
\begin{jieda}
\begin{align*}
\LHS & = \dfrac{\sin(\ang{120}+x)+\sin(\ang{120}-x)}{\sin(\ang{120}+x)\sin(\ang{120}-x)}\\
& = \dfrac{\sqrt{3}\cos{x}}{\sin^2\ang{120}-\sin^2x}\\
& = \dfrac{\sqrt{3}\cos{x}}{\frac{3}{4}-\sin^2x}
\end{align*}
于是$3-4\sin^2x=3\cos{x}$, 故$4\cos^2x-3\cos{x}-1=0$, 解得$\cos{x}=1$或$-\dfrac{1}{4}$. 
\end{jieda}

\item 若$\cos(\alpha+\beta)\cos(\alpha-\beta)=\dfrac{1}{3}$, 则$\cos^2\alpha-\sin^2\beta=$\dotfill(\kh{C})
\xx{$-\dfrac{2}{3}$}{$-\dfrac{1}{3}$}{$\dfrac{1}{3}$}{$\dfrac{2}{3}$}
\begin{jieda}
$\cos(\alpha+\beta)\cos(\alpha-\beta)=\cos^2\alpha-\sin^2\beta$. 
\end{jieda}

\item 若$\alpha,\beta$为锐角, 则下列各式中, 正确的是\dotfill(\kh{C})
\xx{$\cos(\alpha+\beta)>\cos\alpha+\cos\beta$}
{$\cos(\alpha+\beta)>\sin\alpha+\sin\beta$}
{$\cos(\alpha+\beta)<\cos\alpha+\cos\beta$}
{$\cos(\alpha+\beta)<\sin\alpha+\sin\beta$}
\begin{jieda}
令$\alpha=\beta=\frac{\pi}{3}$, 可以否定AB选项; \\
令$\alpha=\beta\to0$, 可以否定D选项; \\
对于C选项, 注意到$\alpha,\beta$为锐角, 则$\cos(\alpha+\beta)<\cos\alpha<\cos\alpha+\cos\beta$. 
\end{jieda}


\item 已知$\sin\alpha+\sin\beta=1$, 则$\cos\alpha+\cos\beta$的取值范围是\dotfill(\kh{C})
\xx{$[0,\sqrt{3}]$}{$[0,2]$}{$[-\sqrt{3},\sqrt{3}]$}{$[-2,2]$}
\begin{jieda}
令$t=\cos\alpha+\cos\beta$, 则$t^2+1=2+2\cos(\alpha-\beta)$, 故$t^2\in[-1,3]$, 故$t\in[-\sqrt{3},\sqrt{3}]$. \\
当$\alpha=\beta=\dfrac{\pi}{6}$时, $t=\sqrt{3}$; 当$\alpha=\beta=\dfrac{5\pi}{6}$时, $t=-\sqrt{3}$. \\
综上, $t\in[-\sqrt{3},\sqrt{3}]$. 
\end{jieda}


\item 设函数$f(x)=\dfrac{\sin3x\sin^3x+\cos3x\cos^3x}{\cos^22x}+\sin2x$, $x\in\left[-\dfrac{\pi}{4},\dfrac{\pi}{4}\right]$的最小值为\dotfill(\kh{D})
\xx{$-1$}{$-\sqrt{2}$}{$1$}{不存在}
\begin{jieda}
这个定义域莫不是有问题? 
\begin{align*}
f(x)
& = \dfrac{\sin(2x+x)\sin^3x+\cos(2x+x)\cos^3x}{\cos^22x}+\sin2x\\
& = \dfrac{\sin2x\cos{x}\sin{x}(\sin^2x-\cos^2x)+\cos2x(1-2\sin^2x\cos^2x)}{\cos^22x}+\sin2x\\
& = \cos2x+\sin2x\\
& = \sqrt{2}\sin\left(2x+\dfrac{\pi}{4}\right)
\end{align*}
\begin{align*}
f(x)
& = \dfrac{\sin3x\sin{x}(1-\cos^2x)+\cos3x\cos{x}(1-\sin^2x)}{\cos^22x}+\sin2x\\
& = \dfrac{(\sin3x\sin{x}+\cos3x\cos{x})-\sin{x}\cos{x}(\sin3x\cos{x}+\cos3x\sin{x})}{\cos^22x}+\sin2x\\
& = \dfrac{\cos2x-\sin{x}\cos{x}\sin4x}{\cos^22x}+\sin2x\\
& = \dfrac{\cos2x-\sin^22x\cos2x}{\cos^22x}+\sin2x\\
& = \cos2x+\sin2x
\end{align*}
因为$x\in\left[-\dfrac{\pi}{4},\dfrac{\pi}{4}\right]$, 故$2x+\dfrac{\pi}{4}\in\left[-\dfrac{\pi}{4},\dfrac{3\pi}{4}\right]$, 
于是$f(x)\in[-1,\sqrt{2}]$. \\
但是定义域呢? 
\end{jieda}


\item 已知$\cos\alpha+\cos\beta=\cos\alpha\cos\beta$, $\cos\dfrac{\alpha+\beta}{2}=\dfrac{4}{5}$, 求$\cos\dfrac{\alpha-\beta}{2}$的值. 
\begin{jieda}
因为$\cos\dfrac{\alpha+\beta}{2}=\dfrac{4}{5}$, 则$\cos(\alpha+\beta)=\dfrac{7}{25}$. \\
$\cos\dfrac{\alpha+\beta}{2}=\dfrac{4}{5}\iff 2\cos\dfrac{\alpha+\beta}{2}\cos\dfrac{\alpha-\beta}{2}=\dfrac{1}{2}(\cos(\alpha+\beta)+\cos(\alpha-\beta))$, \\
故$2\cdot\dfrac{4}{5}\cdot\cos\dfrac{\alpha-\beta}{2}=\dfrac{1}{2}\left(\dfrac{7}{25}+\cos(\alpha-\beta)\right)$, \\
故$\dfrac{16}{5}t=\dfrac{7}{25}+2t^2-1$, 其中$t=\cos\dfrac{\alpha-\beta}{2}$. \\
解得$t=\frac{9}{5}$(舍去)或$-\dfrac{1}{5}$. \\
故$\cos\dfrac{\alpha-\beta}{2}=-\dfrac{1}{5}$. 
\end{jieda}


\item 在$\triangle{ABC}$中, 
\begin{enumerate}[(1)]
\item 已知$\sin^2\dfrac{A}{2}+\sin^2\dfrac{B}{2}+\sin^2\dfrac{C}{2}=\cos^2\dfrac{B}{2}$, 求$\tan\dfrac{A}{2}\tan\dfrac{C}{2}$的值; 
\item 求证: $\sin{A}+\sin{B}+\sin{C}=4\cos\dfrac{A}{2}\cos\dfrac{B}{2}\cos\dfrac{C}{2}$. 
\end{enumerate}
\begin{jieda}
\begin{enumerate}[(1)]
\item 法一: 因为$\sin^2\dfrac{A}{2}+\sin^2\dfrac{B}{2}+\sin^2\dfrac{C}{2}=1-2\sin\dfrac{A}{2}\sin\dfrac{B}{2}\sin\dfrac{C}{2}$. (需要证一下)\\
于是$1-2\sin\dfrac{A}{2}\sin\dfrac{B}{2}\sin\dfrac{C}{2}=\cos^2\dfrac{B}{2}$, 故$\sin^2\dfrac{B}{2}=2\sin\dfrac{A}{2}\sin\dfrac{B}{2}\sin\dfrac{C}{2}$. \\
即$\sin\dfrac{B}{2}=2\sin\dfrac{A}{2}\sin\dfrac{C}{2}$, 故$\sin\dfrac{\pi-A-C}{2}=2\sin\dfrac{A}{2}\sin\dfrac{C}{2}$. \\
于是$\cos\dfrac{A}{2}\cos\dfrac{C}{2}-\sin\dfrac{A}{2}\sin\dfrac{C}{2}=2\sin\dfrac{A}{2}\sin\dfrac{C}{2}$. \\
易得$\tan\dfrac{A}{2}\tan\dfrac{C}{2}=\dfrac{1}{3}$. \\
法二: $\iff\sin^2\dfrac{A}{2}+\sin^2\dfrac{C}{2}=\cos^2\dfrac{B}{2}-\sin^2\dfrac{B}{2}$\\
$\iff\dfrac{1-\cos{A}}{2}+\dfrac{1-\cos{C}}{2}=\cos{B}$\\
$\iff 2+2\cos(A+C)=\cos{A}+\cos{C}$\\
$\iff 4\cos^2\dfrac{A+C}{2}=2\cos\dfrac{A+C}{2}\cos\dfrac{A-C}{2}$\\
$\iff 2\cos\dfrac{A+C}{2}=\cos\dfrac{A-C}{2}$\\
$\iff \cos\dfrac{A}{2}\cos\dfrac{C}{2}=3\sin\dfrac{A}{2}\sin\dfrac{C}{2}$

\item 
\begin{align*}
\LHS & = 2\sin\dfrac{A+B}{2}\cos\dfrac{A-B}{2}+2\sin\dfrac{C}{2}\cos\dfrac{C}{2}\\
& = 2\cos\dfrac{C}{2}\left(\cos\dfrac{A-B}{2}+\sin\dfrac{C}{2}\right)\\
& = 2\cos\dfrac{C}{2}\left(\cos\dfrac{A-B}{2}+\cos\dfrac{A+B}{2}\right)\\
& = 4\cos\dfrac{A}{2}\cos\dfrac{B}{2}\cos\dfrac{C}{2}
\end{align*}
\end{enumerate}
\ps 展一(去掉一角)留二(保留两角)是基本思路. 
\end{jieda}


\item 已知$\sin{x}+\sin{y}=\sqrt{2}$, $\cos{x}+\cos{y}=\dfrac{2\sqrt{3}}{3}$, 求$\cos(x+y)$及$\tan{x}+\tan{y}$的值. 
\begin{jieda}
两式两边平方再相加可得$\cos(x-y)=\frac{2}{3}$. \\
又$\begin{cases}
2\sin\frac{x+y}{2}\cos\frac{x-y}{2}=\sqrt{2}\\
2\cos\frac{x+y}{2}\cos\frac{x-y}{2}=\frac{2\sqrt{3}}{3}
\end{cases}$, 两式相除, 得$\tan\frac{x+y}{2}=\sqrt{\frac{3}{2}}$. \\
于是由万能公式有$\cos(x+y)=-\frac{1}{5}$, $\sin(x+y)=\frac{2\sqrt{6}}{5}$. 
\begin{align*}
\tan{x}+\tan{y} & = \dfrac{\sin{x}}{\cos{x}}+\dfrac{\sin{y}}{\cos{y}}\\
& = \dfrac{\sin(x+y)}{\cos{x}\cos{y}}\\
& = \dfrac{\sin(x+y)}{\frac{1}{2}(\cos(x+y)+\cos(x-y))}\\
& = \dfrac{12\sqrt{6}}{7}
\end{align*}
\ps 另外, 两式两边平方再相减, 得
\begin{align*}
\frac{4}{3}-2
& = \cos^2x+\cos^2y+2\cos{x}\cos{y}-(\sin^2x+\sin^2y+2\sin{x}\sin{y})\\
& = \cos2x+\cos2y+2\cos(x+y)\\
& = 2\cos(x+y)\cos(x-y)+2\cos(x+y)
\end{align*}
于是在$\cos(x-y)=\frac{2}{3}$的前提下, 亦可得到$\cos(x+y)=-\frac{1}{5}$. 
\end{jieda}


\item 设$\alpha,\beta,\gamma$均为锐角, 则等式$\cos^2\alpha+\cos^2\beta+\cos^2\gamma+2\cos\alpha\cos\beta\cos\gamma=1$成立的充分必要条件是$\alpha+\beta+\gamma=\pi$. 
\begin{jieda}
法一: $\iff\dfrac{1+\cos2\alpha}{2}+\dfrac{1+\cos2\beta}{2}+\cos^2\gamma+2\cos\alpha\cos\beta\cos\gamma=1$\\
$\iff\dfrac{\cos2\alpha+\cos2\beta}{2}+\cos^2\gamma+2\cos\alpha\cos\beta\cos\gamma=0$\\
$\iff\cos(\alpha+\beta)\cos(\alpha-\beta)+\cos^2\gamma+(\cos(\alpha+\beta)+\cos(\alpha-\beta))\cos\gamma=0$\\
$\iff(\cos\gamma+\cos(\alpha+\beta))(\cos\gamma+\cos(\alpha-\beta))=0$\\
因为$\alpha,\beta,\gamma$均为锐角, 故$\cos\gamma>0$, $\cos(\alpha-\beta)>0$, \\
于是$\cos\gamma+\cos(\alpha+\beta)=0$, 故$\gamma=\pi-\alpha-\beta$, 即$\alpha+\beta+\gamma=\pi$. \\
法二: $\iff(\cos\alpha+\cos\beta\cos\gamma)^2=\sin^2\beta\sin^2\gamma$. \\
由于$\alpha,\beta,\gamma$均为锐角, 故$\cos\alpha+\cos\beta\cos\gamma=\sin\beta\sin\gamma$, 即$\cos\alpha+\cos(\beta+\gamma)=0$. \\
法三: 注意到$\alpha,\beta,\gamma$均为锐角, 则$\cos\gamma=-\cos\alpha\cos\beta+\sqrt{(1-\cos^2\alpha)(1-\cos^2\beta)}$. \\
于是$\cos\gamma=-\cos\alpha\cos\beta+\sin\alpha\sin\beta=-\cos(\alpha+\beta)$. 
\end{jieda}

\end{enumerate}





\section{解三角形}


\subsection{解三角形(1)}
\zsd 正弦定理


\begin{enumerate}
\item  在$\triangle{ABC}$中, $a=5$, $B=\ang{105}$, $C=\ang{15}$, 则此三角形的最大边的长为\blkda[3]{$\dfrac{5\sqrt{6}+15\sqrt{2}}{6}$}.
\begin{jieda}
$A=\ang{60}$, $b=\dfrac{a}{\sin{A}}\sin{B}$. 
\end{jieda}

\item 海上有$A$、$B$、$C$三个小岛, 测得$A$、$B$两岛相距10海里, $\angle{BAC}=\ang{60}$, $\angle{ABC}=\ang{75}$, 则$B$、$C$两岛间的距离是\blkda[2]{$5\sqrt{6}$海里}. 
\begin{jieda}
$\dfrac{AB}{\sin\angle{ACB}}=\dfrac{BC}{\sin\angle{BAC}}$. 
\end{jieda}

\item 一船自西向东匀速航行, 上午10时到达一座灯塔$P$的南偏西$\ang{75}$距塔68海里的$M$处, 下午2时到达这座灯塔的东南方向的$N$处, 
则这只船的航行速度为\blkda{$\dfrac{17\sqrt{6}}{2}$} 海里/小时. 
\begin{jieda}
$\dfrac{68}{\sin\ang{45}}=\dfrac{MN}{\sin\ang{120}}$, 故$MN=34\sqrt{6}$, $v=\dfrac{MN}{4}$. 
\end{jieda}

\item 在$\triangle{ABC}$中, 已知$\cos{A}=\dfrac{5}{13}$, $\sin{B}=\dfrac{3}{5}$, 则$\cos{C}$的值为\blkda{$\dfrac{16}{65}$}. 
\begin{jieda}
$\sin{A}=\dfrac{12}{13}$, 则$\sin{A}>\sin{B}$, 故$A>B$. \\
又$A$为锐角, 故$B$亦为锐角, 于是$\cos{B}=\dfrac{4}{5}$. 
\end{jieda}

\item 在$\triangle{ABC}$中, 若$b=2a\sin{B}$, 则这个三角形中角$A$的值是\blkda{$\dfrac{\pi}{6}$或$\dfrac{5\pi}{6}$}. 
\begin{jieda}
$\iff \sin{B}=2\sin{A}\sin{B}\iff \sin{A}=\frac{1}{2}$. 
\end{jieda}

\item 在$\triangle{ABC}$中, 若$\tan{A}=\dfrac{1}{3}$, $C=\ang{105}$, $BC=1$, 则$AB$=\blkda{$\dfrac{\sqrt{10}}{2}$}. 
\begin{jieda}
$\dfrac{BC}{\sin{A}}=\dfrac{AB}{\sin{C}}$. 
\end{jieda}

\item 在$\triangle{ABC}$中, 若$A=2B$, 则$a$等于\dotfill(\kh{D})
\xx{$2b\sin{A}$}{$2b\cos{A}$}{$2b\sin{B}$}{$2b\cos{B}$}
\begin{jieda}
$\dfrac{a}{\sin{A}}=\dfrac{b}{\sin{B}}$, 则$\dfrac{a}{\sin2B}=\dfrac{b}{\sin{B}}$, 故$a=2b\cos{B}$. 
\end{jieda}

\item 在$\triangle{ABC}$中, $a=6$, $B=\ang{30}$, $C=\ang{120}$, 则$\triangle{ABC}$的面积是\dotfill(\kh{C})
\xx{9}{18}{$9\sqrt{3}$}{$18\sqrt{3}$}
\begin{jieda}
$S=2R^2\sin{A}\sin{B}\sin{C}=\dfrac{1}{2}\left(\dfrac{a}{\sin{A}}\right)^2\sin{A}\sin{B}\sin{C}$. 
\end{jieda}

\item 在$\triangle{ABC}$中, 若$2\cos{B}\sin{A}=\sin{C}$, 则$\triangle{ABC}$的形状一定是\dotfill(\kh{C})
\xx{等腰直角三角形}{直角三角形}{等腰三角形}{等边三角形}
\begin{jieda}
$\iff 2\cos{B}a=c\iff 2a\cos{B}=a\cos{B}+b\cos{A}\iff a\cos{B}=b\cos{A}\iff \sin{A}\cos{B}=\sin{B}\cos{A}\iff\sin(A-B)=0$. 
\end{jieda}

\item 在$\triangle{ABC}$中, 由已知条件解三角形, 其中有两解的是\dotfill(\kh{C})
\xx{$b=20$, $A=\ang{45}$, $C=\ang{80}$}
{$a=30$, $c=28$, $B=\ang{60}$}
{$a=14$, $b=16$, $A=\ang{45}$}
{$a=12$, $c=15$, $A=\ang{120}$}
\begin{jieda}
三角形全等, 画图. 
\end{jieda}

\item 在$\triangle{ABC}$中, 已知$a=20\si{cm}$, $b=28\si{cm}$, $A=\ang{40}$, 解三角形(角度精确到$\ang{1}$, 边长精确到1\si{cm}). 
\begin{jieda}
$(B,C,c)=(\ang{64},\ang{76},30\si{cm}), (\ang{116},\ang{24},13\si{cm})$. 
\end{jieda}

\item 在$\triangle{ABC}$中, $\sin{A}+\cos{A}=\dfrac{\sqrt{2}}{2}$, $AC=2$, $AB=3$, 求$\tan{A}$的值和$\triangle{ABC}$的面积. 
\begin{jieda}

$\tan{A}=-2-\sqrt{3}$, $S=\dfrac{3\sqrt{6}+3\sqrt{2}}{4}$. 
\end{jieda}

\item 在$\triangle{ABC}$中, $A$、$B$为锐角, 角$A$、$B$、$C$所对的边分别为$a$、$b$、$c$, 且$\sin{A}=\dfrac{\sqrt{5}}{5}$, $\sin{B}=\dfrac{\sqrt{10}}{10}$. 
\begin{enumerate}
\item 求$A+B$的值; 
\item 若$a-b=\sqrt{2}-1$, 求$a$、$b$、$c$的值. 
\end{enumerate}
\begin{jieda}
\begin{enumerate}
\item $\cos(A+B)=\dfrac{\sqrt{2}}{2}$, 故$A+B=\dfrac{\pi}{4}$. 
\item $\dfrac{a}{b}=\dfrac{\sin{A}}{\sin{B}}=\sqrt{2}$, 结合$a-b=\sqrt{2}-1$, 可得$a=\sqrt{2}$, $b=1$. \\
正弦定理: $c=\dfrac{a}{\sin{A}}\sin{C}=\sqrt{5}$. \\
余弦定理: $c^2=a^2+b^2-2ab\cos{C}=a^2+b^2+2ab\cos(A+B)$, 故$c=\sqrt{5}$. 
\end{enumerate}
\end{jieda}

\item 在$\triangle{ABC}$中, 记角$A$、$B$、$C$所对的边分别为$a$、$b$、$c$, 已知$2\sin^2\dfrac{A+B}{2}=1+\cos2C$. 
\begin{enumerate}
\item 求$C$的大小; 
\item 若$c^2=2b^2-2a^2$, 求$\cos2A-\cos2B$的值. 
\end{enumerate}
\begin{jieda}
$\iff 0=\cos(A+B)+\cos2C=-\cos{C}+2\cos^2C-1$, 解得$\cos{C}=-\dfrac{1}{2}$或1(舍去). 
\begin{enumerate}
\item 因为$C\in(0,\pi)$, 故$C=\dfrac{2\pi}{3}$; 
\item $\iff \sin^2C=2\sin^2B-2\sin^2A\iff \dfrac{3}{4}=(1-\cos2B)-(1-\cos2A)$, \\
故$\cos2A-\cos2B=\dfrac{3}{4}$. 
\end{enumerate}
\end{jieda}

\end{enumerate}



% 缺图 20240511
\subsection{解三角形(2)}
\zsd 正余弦定理

\begin{enumerate}
\item 在$\triangle{ABC}$中, $(2a-c)\cos{B}=b\cos{C}$, 则$B$=\blkda{$\dfrac{\pi}{3}$}. 
\begin{jieda}
$\iff 2a\cos{B}=c\cos{B}+b\cos{C}=a\iff\cos{B}=\frac{1}{2}$. 
\end{jieda}

\item 在$\triangle{ABC}$中, 若其面积$S=\dfrac{a^2+b^2-c^2}{4\sqrt 3}$, 则$\angle{C}$=\blkda{$\dfrac{\pi}{6}$}. 
\begin{jieda}
$\iff \dfrac{1}{2}ab\sin{C}=\dfrac{2ab\cos{C}}{4\sqrt{3}}$. 
\end{jieda}

\item 已知$\triangle{ABC}$的面积是30, 且$\cos{A}=\dfrac{12}{13}$, $c-b=1$, 则$a$=\blkda{5}. 
\begin{jieda}
$30=\dfrac{1}{2}bc\sin{A}$, 则$bc=13\cdot 12$. 
\end{jieda}

\item 在$\triangle{ABC}$中, $(a+b+c)(\sin{A}+\sin{B}-\sin{C})=3a\sin{B}$, 则$\angle{C}$=\blkda{$\dfrac{\pi}{3}$}. 
\begin{jieda}
$\iff(a+b)^2-c^2=3ab\iff a^2+b^2-c^2=ab\iff 2ab\cos{C}=ab$. 
\end{jieda}

\item 在$\triangle{ABC}$中, $AB=2$, $AC=3$, $A=\dfrac{\pi}{3}$, 则$\triangle{ABC}$外接圆的半径为\blkda{$\dfrac{\sqrt{21}}{3}$}. 
\begin{jieda}
$2R=\dfrac{a}{\sin{A}}$, $a^2=b^2+c^2-2bc\cos{A}$. 
\end{jieda}

\item 在锐角$\triangle{ABC}$中, $\dfrac{b}{a}+\dfrac{a}{b}=6\cos{C}$, 则$\dfrac{\tan{C}}{\tan{A}}+\dfrac{\tan{C}}{\tan{B}}$=\blkda{4}. 
\begin{jieda}
$\iff a^2+b^2=6ab\cos{C}=3(a^2+b^2-c^2)\iff3c^2=2(a^2+b^2)$, 
于是
\begin{align*}
\dfrac{\tan{C}}{\tan{A}}+\dfrac{\tan{C}}{\tan{B}}
& =\dfrac{\sin{C}}{\cos{C}}\left(\dfrac{\cos{A}}{\sin{A}}+\dfrac{\cos{B}}{\sin{B}}\right)\\
& =\dfrac{\sin{C}}{\cos{C}}\cdot\dfrac{\cos{A}\sin{B}+\cos{B}\sin{A}}{\sin{A}\sin{B}}\\
& =\dfrac{\sin{C}}{\cos{C}}\cdot\dfrac{\sin(A+B)}{\sin{A}\sin{B}}\\
& =\dfrac{c^2}{ab\cos{C}}\\
& =\dfrac{2c^2}{a^2+b^2-c^2}\\
& =4.
\end{align*}
\end{jieda}

\item 在$\triangle{ABC}$中, $a^2-b^2=\sqrt{3}bc$, $\sin{C}=2\sqrt{3}\sin{B}$, 则$A=$\dotfill(\kh{A})
\xx{$\ang{30}$}{$\ang{60}$}{$\ang{120}$}{$\ang{150}$}
\begin{jieda}
$c=2\sqrt{3}b$, 于是$a^2-b^2=\sqrt{3}b\cdot2\sqrt{3}b=6b^2$, 即$a=\sqrt{7}b$. 故$a:b:c=\sqrt{7}:1:2\sqrt{3}$. 
\end{jieda}

\item 某人在$C$点测得某塔在南偏西$\ang{80}$, 塔顶仰角为$\ang{45}$, 此人沿南偏东$\ang{40}$方向前进10米到$D$, 测得塔顶$A$的仰角为$\ang{30}$, 则塔高为\dotfill(\kh{C})
\xx{15米}{5米}{10米}{12米}
\begin{jieda}
$3h^2=h^2+100-2h\cdot10\cdot\cos\ang{120}$. 
\end{jieda}

\item 若满足条件$C=\dfrac{\pi}{3}$, $AB=\sqrt{3}$, $BC=a$的$\triangle{ABC}$有两个, 那么$a$的取值范围是\dotfill(\kh{C})
\xx{$(1,\sqrt{2})$}{$(\sqrt{2},\sqrt{3})$}{$(\sqrt{3},2)$}{$(1,2)$}
\begin{jieda}
法一: 画图. \\
法二: 余弦定理. 设$AC=x$, 则$3=x^2+a^2-2ax\cos\ang{60}$. \\
于是$x^2-ax+a^2-3=0$, 该方程有两个不等正根. 
\end{jieda}

\item 在$\triangle{ABC}$中, $(a^2+c^2-b^2)\tan{B}=\sqrt{3}ac$, 则角$B$的值为\dotfill(\kh{D})
\xx{$\dfrac{\pi}{6}$}
{$\dfrac{\pi}{3}$}
{$\dfrac{\pi}{6}$或$\dfrac{5\pi}{6}$}
{$\dfrac{\pi}{3}$或$\dfrac{2\pi}{3}$}
\begin{jieda}
$\sin{B}=\frac{\sqrt{3}}{2}$. 
\end{jieda}

\item 如图, 一货轮航行到$M$处, 测得灯塔$S$在货轮的北偏东$\ang{15}$, 与灯塔$S$相距20海里, 随后货轮按北偏西$\ang{30}$的方向航行30分钟后, 
又测得灯塔在货轮的东北方向, 则货轮的速度为\dotfill(\kh{B})
\xx
{$20(\sqrt{2}+\sqrt{6})$海里/小时}{$20(\sqrt{6}-\sqrt{2})$海里/小时}
{$20(\sqrt{6}+\sqrt{3})$海里/小时}{$20(\sqrt{6}-\sqrt{3})$海里/小时}

\item 在锐角$\triangle{ABC}$中, $a=2b\sin{A}$.
\begin{enumerate}
\item 求$B$的大小; 
\item 若$a=3\sqrt{3}$, $c=5$, 求$b$. 
\end{enumerate}
\begin{jieda}
\begin{enumerate}
\item $B=\dfrac{\pi}{6}$; 
\item $b=\sqrt{7}$. 
\end{enumerate}
\end{jieda}

\item 在$\triangle{ABC}$中, 若$b^2=ac$, 且$a^2-c^2=ac-bc$, 求$\angle{A}$的大小及$\dfrac{b\sin{B}}{c}$的值. 
\begin{jieda}
$\cos{A}=\frac{1}{2}$, 故$A=\dfrac{\pi}{3}$, $\dfrac{b\sin{B}}{c}=\dfrac{\sqrt{3}}{2}$. 
\end{jieda}

\item 在$\triangle{ABC}$中, $2a\sin{A}=(2b+c)\sin{B}+(2c+b)\sin{C}$. 
\begin{enumerate}
\item 求$A$的大小; 
\item 若$\sin{B}+\sin{C}=1$, 试判断$\triangle{ABC}$的形状. 
\end{enumerate}
\begin{jieda}
\begin{enumerate}
\item $\cos{A}=-\frac{1}{2}$, 故$A=\frac{2\pi}{3}$. 
\item $1=2\sin\frac{B+C}{2}\cos\frac{B-C}{2}$, $B+C=\frac{\pi}{3}$, 故$\cos\frac{B-C}{2}=1$. \\
易得$B=C$, 故等腰. 
\end{enumerate}
\end{jieda}

\item 在$\triangle{ABC}$中, 若$(a+b+c)(a-b+c)=3ac$, 且$\tan{A}+\tan{C}=3+\sqrt{3}$, $AB$边上的高为$4\sqrt{3}$, 求角$A$、$B$、$C$的大小与边$a$、$b$、$c$的长.
\begin{jieda}
$a^2+c^2-b^2=ac$, 故$\cos{B}=\frac{1}{2}$, 即$B=\frac{\pi}{3}$. \\
于是$A+C=\frac{2\pi}{3}$, 又$\tan{A}+\tan{C}=3+\sqrt{3}$, \\
故$\tan{A}=1$或$2+\sqrt{3}$. \\
\ps 可以$\tan{B}=-\tan(A+C)$, 亦可$\tan{A}+\tan{C}=\dfrac{\sin(A+C)}{\cos{A}\cos{C}}$. \\
$(A,B,C,a,b,c)=(\frac{\pi}{4},\frac{\pi}{3},\frac{5\pi}{12},8,4\sqrt{6},4+4\sqrt{3}), (\frac{5\pi}{12},\frac{\pi}{3},\frac{\pi}{4},8,12\sqrt{2}-4\sqrt{6},8\sqrt{3}-8)$. 
\end{jieda}

\end{enumerate}



% 缺图. 20240511
\subsection{解三角形(3)}
\zsd 正余弦定理

\begin{enumerate}
\item 在$\triangle{ABC}$中, $(\sqrt{3}b-c)\cos{A}=a\cos{C}$, 则$\cos{A}$=\blkda{$\dfrac{\sqrt{3}}{3}$}. 
\begin{jieda}
$b=c\cos{A}+a\cos{C}$. 
\end{jieda}

\item 在$\triangle{ABC}$中, $A=\dfrac{\pi}{6}$, $(1+\sqrt{3})c=2b$, 则$C=$\blkda{$\dfrac{\pi}{4}$}. 

\item 在$\triangle{ABC}$中, $b^2+c^2=a^2+\sqrt{2}bc$, 且$a=\sqrt{2}b$, 则$\angle{C}$=\blkda{$\dfrac{7\pi}{12}$}. 

\item 在$\triangle{ABC}$中, $\cos\dfrac{A}{2}=\dfrac{2\sqrt{5}}{5}$, $\triangle{ABC}$的面积为2, 且$b+c=6$, 则$a$=\blkda{$2\sqrt{5}$}. 

\item 在锐角$\triangle{ABC}$中, $a=1$, $b=2$, 则边长$c$的取值范围是\blkda{$(\sqrt{3},\sqrt{5})$}. 
\begin{jieda}
三个角都是锐角, 又$a<b$, 故只要$c^2>a^2+b^2\land b^2>c^2+a^2$. 当然要检验``两边之和大于第三边''. 
\end{jieda}

\item 已知$\triangle{ABC}$的周长为$\sqrt{2}+1$, 且$\sin{A}+\sin{B}=\sqrt{2}\sin{C}$. 若$\triangle{ABC}$的面积为$\dfrac{1}{6}\sin{C}$, 则$C$=\blkda{$\dfrac{\pi}{3}$}. 
\begin{jieda}

\end{jieda}

\item 在$\triangle{ABC}$中, $A=\ang{60}$, $b=1$, $S_{\triangle{ABC}}=\sqrt{3}$, 则$\dfrac{a+b+c}{\sin{A}+\sin{B}+\sin{C}}$=\blkda{$\dfrac{2\sqrt{39}}{3}$}. 
\begin{jieda}

\end{jieda}

\item 在$\triangle{ABC}$中, ``$c\cos{B}=b\cos{C}$''是``$\triangle{ABC}$是等腰三角形''的\dotfill(\kh{A})
\xx{充分不必要条件}{必要不充分条件}
{充分必要条件}{既不充分也不必要条件}
\begin{jieda}
条件$\iff c\cdot\dfrac{a^2+c^2-b^2}{2ac}=b\cdot\dfrac{a^2+b^2-c^2}{2ab}\iff a^2+c^2-b^2=a^2+b^2-c^2\iff b=c$. \\
等腰$\iff a=b\lor b=c\lor c=a$. 
\end{jieda}

\item 若$\triangle{ABC}$的三个内角满足$\sin{A}:\sin{B}:\sin{C}=5:11:13$, 则$\triangle{ABC}$\dotfill(\kh{C})
\xx{一定是锐角三角形}{一定是直角三角形}
{一定是钝角三角形}{可能是锐角三角形, 也可能是钝角三角形}
\begin{jieda}
$a:b:c=5:11:13$, 则$\cos{C}<0$. 
\end{jieda}


\item $\triangle{ABC}$中, $2b=a+c$, $\angle{B}=\ang{30}$, 且$\triangle{ABC}$的面积为$\dfrac{3}{2}$, 则$b=$\dotfill(\kh{B}) 
\xx{$\dfrac{1+\sqrt{3}}{2}$}
{$1+\sqrt{3}$}
{$\dfrac{2+\sqrt{3}}{2}$}
{$2+\sqrt{3}$}
\begin{jieda}
$(a+c)^2=4b^2=4(a^2+c^2-2ac\cos{B})$, $S=\frac{1}{2}ac\sin{B}$. 
\end{jieda}

\item 在锐角$\triangle{ABC}$中, $\sqrt{3}a=2c\sin{A}$. 
\begin{enumerate}
\item 求角$C$的大小;
\item 若$c=\sqrt{7}$, 且$\triangle{ABC}$的面积为$\dfrac{3\sqrt{3}}{2}$, 求$a+b$的值. 
\end{enumerate}
\begin{jieda}
\begin{enumerate}
\item $\dfrac{\pi}{3}$

\item $a+b=5$
\end{enumerate}
\end{jieda}


\item 已知半径为$R$的圆外接于$\triangle{ABC}$, 且$2R(\sin^2A-\sin^2C)=(\sqrt{3}a-b)\sin{B}$. 
\begin{enumerate}
\item 求角$C$; 
\item 求$\triangle{ABC}$面积的最大值. 
\end{enumerate}
\begin{jieda}
\begin{enumerate}
\item $\iff a^2-c^2=(\sqrt{3}a-b)b\iff a^2+b^2-c^2=\sqrt{3}ab$, 于是$2ab\cos{C}=\sqrt{3}ab$. \\
于是$\cos{C}=\dfrac{\sqrt{3}}{2}$, 故$C=\dfrac{\pi}{6}$. 

\item 因为面积必须是长度的平方, 条件中又没有边长, 故$S$用$R$表示. \\
$S=2R^2\sin{A}\sin{B}\sin{C}=R^2\sin{C}(\cos(A-B)-\cos(A+B))=\frac{1}{2}R^2(\cos(A-B)+\frac{\sqrt{3}}{2})$, \\
故$S_{\max}=\dfrac{2+\sqrt{3}}{4}R^2(A=B=\frac{5\pi}{12})$
\end{enumerate}
\end{jieda}

\item 如图, 半圆$O$的直径为2, $A$为直径延长线上的一点, $OA=2$, $B$为半圆上任意一点, 以$AB$为一边作等边$\triangle{ABC}$. 
问: 点$B$在什么位置时, 四边形$OACB$面积最大? 
\begin{jieda}
设$\angle{AOB}=\alpha$, 则$S=2\sin\left(\alpha-\dfrac{\pi}{3}\right)+\dfrac{5\sqrt{3}}{4}$. \\
于是当$\alpha=\dfrac{5\pi}{6}$时, $S_{\max}=2+\dfrac{5\sqrt{3}}{4}$. 
\end{jieda}

\item 在一个特定时段内, 以点$E$为中心的7海里以内海域被设为警戒水域, 点$E$正北55海里处有一个雷达观测站$A$. 
某时刻测得一艘匀速直线行驶的船只位于点$A$北偏东$\ang{45}$且与点$A$相距$40\sqrt{2}$海里的位置$B$, 经过40分钟又测得该船已行驶到点$A$北偏东$\ang{45}+\theta$
(其中$\sin\theta=\frac{\sqrt{26}}{26}$, $\ang{0}<\theta<\ang{90}$)且与点$A$相距$10\sqrt{13}$海里的位置$C$. 
\begin{enumerate}
\item 求该船的行驶速度(单位: 海里/小时); 
\item 若该船不改变航行方向继续行驶, 判断它是否会进入警戒水域, 并说明理由. 
\end{enumerate}
\begin{jieda}
\begin{enumerate}
\item $v=15\sqrt{5}$海里/小时

\item $d=3\sqrt{5}<7$, 故会进入警戒水域. 
\end{enumerate}
\end{jieda}

\end{enumerate}

